% ===========================================================================
%  Projet GANDAL, equipe 20
%  Tache T005 : mise en place du service DNS
%  Auteur : FOMETHE SOMBANANG Maximilien
%  Compilation : pdflatex main.tex  (deux passes pour la table des matieres)
% ===========================================================================
\documentclass[11pt,a4paper,oneside]{report}
% ---------------------------------------------------------------------------
%  Preambule : typographie Latin Modern, mise en page et environnements
%  Projet GANDAL, tache T005
% ---------------------------------------------------------------------------

\usepackage[T1]{fontenc}
\usepackage[utf8]{inputenc}
\usepackage{lmodern}          % Latin Modern : Computer Modern au format moderne
\usepackage[french]{babel}
% babel-french emploie par defaut un tiret long comme puce de liste :
% on retablit les labels standard, le tiret cadratin etant proscrit.
\frenchsetup{StandardLists=true}
% thebibliography compose son propre titre de chapitre : on le renomme
% plutot que d'en ajouter un second.
\addto\captionsfrench{\renewcommand{\bibname}{Références}}
\usepackage{microtype}        % protrusion et expansion : gris typographique regulier
\usepackage{amsmath,amssymb}
\usepackage{textcomp}

\usepackage[a4paper,top=27mm,bottom=25mm,left=26mm,right=24mm,
            headheight=14pt,headsep=10mm,footskip=13mm]{geometry}

\usepackage{xcolor}
\definecolor{gdnavy}{HTML}{1B3A57}
\definecolor{gdteal}{HTML}{146B6B}
\definecolor{gdgold}{HTML}{B08432}
\definecolor{gdgrey}{HTML}{5C6670}
\definecolor{gdrule}{HTML}{C9CDD2}
\definecolor{gdlight}{HTML}{F2F1EC}
\definecolor{gdcode}{HTML}{F5F6F7}
\definecolor{gdbox}{HTML}{EDF3F3}
\definecolor{gdwarn}{HTML}{FAF5EA}
\definecolor{gdred}{HTML}{A33A3A}

% --------------------------------------------------------------- sectionnement
\usepackage{titlesec}
\titleformat{\chapter}[display]
  {\normalfont\LARGE\bfseries\color{gdnavy}}
  {\normalsize\normalfont\scshape\color{gdteal}\chaptertitlename~\thechapter}
  {6pt}
  {\LARGE}
  [{\vspace{4pt}\color{gdgold}\titlerule[1.1pt]}]
\titlespacing*{\chapter}{0pt}{-16pt}{24pt}

\titleformat{\section}
  {\normalfont\large\bfseries\color{gdnavy}}{\thesection}{0.8em}{}
\titleformat{\subsection}
  {\normalfont\normalsize\bfseries\color{gdteal}}{\thesubsection}{0.7em}{}
\titleformat{\subsubsection}[runin]
  {\normalfont\normalsize\bfseries}{}{0pt}{}[.\,]
\titlespacing*{\section}{0pt}{16pt plus 3pt minus 2pt}{7pt}
\titlespacing*{\subsection}{0pt}{12pt plus 2pt minus 2pt}{5pt}

\setcounter{secnumdepth}{2}
\setcounter{tocdepth}{2}

% ------------------------------------------------------------- en-tetes et pieds
\usepackage{fancyhdr}
\fancypagestyle{gandal}{%
  \fancyhf{}%
  \renewcommand{\headrulewidth}{0.4pt}%
  \renewcommand{\footrulewidth}{0.4pt}%
  \fancyhead[L]{\footnotesize\color{gdgrey}\leftmark}%
  \fancyfoot[L]{\scriptsize\color{gdgrey}\textsc{Projet GANDAL \quad Cellule
                Réseaux et Sécurité \quad Tâche T005}}%
  \fancyfoot[R]{\small\color{gdnavy}\textbf{\thepage}}%
}
\fancypagestyle{plain}{%
  \fancyhf{}%
  \renewcommand{\headrulewidth}{0pt}%
  \renewcommand{\footrulewidth}{0.4pt}%
  \fancyfoot[L]{\scriptsize\color{gdgrey}\textsc{Projet GANDAL \quad Cellule
                Réseaux et Sécurité \quad Tâche T005}}%
  \fancyfoot[R]{\small\color{gdnavy}\textbf{\thepage}}%
}
\renewcommand{\chaptermark}[1]{\markboth{\thechapter\quad #1}{}}
\pagestyle{gandal}

% ------------------------------------------------------------------- tableaux
\usepackage{booktabs}
\usepackage{array}
\usepackage{longtable}
\usepackage{multirow}
\usepackage{colortbl}
\renewcommand{\arraystretch}{1.22}
\setlength{\tabcolsep}{5pt}

% Colonnes de largeur fractionnaire : la somme des fractions vaut 1
\newcolumntype{R}[1]{>{\raggedright\arraybackslash}p{\dimexpr#1\textwidth-2\tabcolsep-0.4pt}}
\newcolumntype{C}[1]{>{\centering\arraybackslash}p{\dimexpr#1\textwidth-2\tabcolsep-0.4pt}}

\newcommand{\thead}[1]{\textbf{\color{white}#1}}
\newcommand{\entete}{\rowcolor{gdnavy}}

% --------------------------------------------------------------------- legendes
\usepackage{caption}
\captionsetup{font=small,labelfont=bf,labelsep=period,width=0.94\textwidth}
\captionsetup[table]{skip=5pt}
\captionsetup[figure]{skip=9pt}

% ----------------------------------------------------------------------- listes
\usepackage{enumitem}
\setlist[itemize]{label=\textbullet,leftmargin=1.35em,topsep=4pt,
                  itemsep=2.5pt,parsep=0pt}
\setlist[itemize,2]{label=\textendash,leftmargin=1.2em}
\setlist[enumerate]{leftmargin=1.75em,topsep=4pt,itemsep=2.5pt,parsep=0pt}

% ------------------------------------------------------------------- listings
\usepackage{listings}
\lstdefinestyle{gandal}{%
  basicstyle=\ttfamily\scriptsize,
  commentstyle=\color{gdgrey}\itshape,
  columns=fullflexible,
  keepspaces=true,
  showstringspaces=false,
  upquote=true,
  breaklines=true,
  breakatwhitespace=false,
  breakindent=1.6em,
  postbreak=\mbox{\textcolor{gdteal}{$\hookrightarrow$}\space},
  frame=leftline,
  framesep=7pt,
  framerule=1.1pt,
  rulecolor=\color{gdteal},
  backgroundcolor=\color{gdcode},
  xleftmargin=10pt,
  xrightmargin=1pt,
  aboveskip=10pt,
  belowskip=4pt,
  literate=%
    {é}{{\'e}}1 {è}{{\`e}}1 {ê}{{\^e}}1 {ë}{{\"e}}1
    {à}{{\`a}}1 {â}{{\^a}}1 {ù}{{\`u}}1 {û}{{\^u}}1
    {î}{{\^i}}1 {ï}{{\"i}}1 {ô}{{\^o}}1 {ç}{{\c c}}1
    {É}{{\'E}}1 {È}{{\`E}}1 {À}{{\`A}}1 {Ç}{{\c C}}1
}
\lstset{style=gandal}
\lstdefinelanguage{pdnsconf}{%
  morecomment=[l]{\#},
  sensitive=true,
}
\lstdefinelanguage{luapdns}{%
  morecomment=[l]{--},
  morestring=[b]",
  morekeywords={local,function,end,if,then,else,return,nil,for,in,do,and,or,not},
  sensitive=true,
}
\renewcommand{\lstlistingname}{Listing}

% ------------------------------------------------------------------- encadres
\usepackage[most]{tcolorbox}
\newtcolorbox{encadre}[1][]{%
  enhanced, breakable, boxrule=0pt, frame hidden,
  borderline west={2.4pt}{0pt}{gdteal},
  colback=gdbox, sharp corners, left=9pt, right=9pt, top=7pt, bottom=7pt,
  fontupper=\small, #1}
\newtcolorbox{encadredecision}[1][]{%
  enhanced, breakable, boxrule=0pt, frame hidden,
  borderline west={2.4pt}{0pt}{gdgold},
  colback=gdwarn, sharp corners, left=9pt, right=9pt, top=7pt, bottom=7pt,
  fontupper=\small, #1}
\newcommand{\titreencadre}[1]{%
  {\footnotesize\bfseries\scshape\color{gdteal}#1\par}\vspace{2pt}}
\newcommand{\titreencadredec}[1]{%
  {\footnotesize\bfseries\scshape\color{gdgold}#1\par}\vspace{2pt}}

% --------------------------------------------------------------------- figures
\usepackage{tikz}
\usetikzlibrary{babel,positioning,arrows.meta,calc,fit,backgrounds,shapes.geometric,
                shadows.blur,decorations.pathreplacing}
\usepackage{float}

\tikzset{
  bloc/.style   ={draw=gdnavy, line width=0.6pt, rounded corners=1.5pt,
                  align=center, inner sep=4pt, font=\scriptsize, fill=white},
  blocteal/.style={bloc, draw=gdteal},
  blocgold/.style={bloc, draw=gdgold},
  blocrouge/.style={bloc, draw=gdred},
  blocgris/.style={bloc, draw=gdgrey},
  hote/.style   ={draw=gdnavy, line width=0.7pt, rounded corners=2pt,
                  fill=gdlight, inner sep=5pt},
  bande/.style  ={draw=gdteal, line width=0.6pt, dash pattern=on 2pt off 2pt,
                  fill=gdbox, rounded corners=2pt, inner sep=6pt},
  fl/.style     ={-{Stealth[length=4pt,width=3pt]}, line width=0.65pt,
                  draw=gdteal},
  flnavy/.style ={fl, draw=gdnavy},
  flgris/.style ={fl, draw=gdgrey, dash pattern=on 2pt off 1.6pt},
  flgold/.style ={fl, draw=gdgold},
  flrouge/.style={fl, draw=gdred, dash pattern=on 2.4pt off 1.8pt},
  etiq/.style   ={font=\tiny, inner sep=1.5pt, text=gdnavy},
  etiqt/.style  ={etiq, text=gdteal},
  etiqr/.style  ={etiq, text=gdred},
  etiqg/.style  ={etiq, text=gdgrey},
  titrepan/.style={font=\scriptsize\bfseries\scshape, text=gdteal, anchor=west},
}

% ------------------------------------------------------------------ references
\usepackage{url}
\usepackage[hidelinks,pdfusetitle]{hyperref}
\hypersetup{
  colorlinks=true, linkcolor=gdnavy, citecolor=gdnavy, urlcolor=gdteal,
  pdftitle={Mise en place du service DNS GANDAL, tache T005},
  pdfauthor={FOMETHE SOMBANANG Maximilien},
  pdfsubject={Projet GANDAL, cellule Reseaux et Securite, tache T005},
  pdfkeywords={GANDAL, DNS, PowerDNS, isolation multi-tenant, TSIG, DNSSEC},
}

% --------------------------------------------------------------------- notations
\newcommand{\zone}[1]{\texttt{#1}}
\newcommand{\dir}[1]{\texttt{#1}}
\newcommand{\opt}[1]{\texttt{#1}}
\newcommand{\ip}[1]{\texttt{#1}}
\newcommand{\statut}[1]{\textsc{#1}}
\newcommand{\exig}[1]{\textsc{#1}}

\newcommand{\interligne}{\vspace{2pt}}
\setlength{\parskip}{3pt plus 1pt minus 0.5pt}
\setlength{\emergencystretch}{2em}
\widowpenalty=10000
\clubpenalty=10000

% siunitx absent de cette installation : passe-plat defensif
\providecommand{\num}[1]{#1}

% ---------------------------------------------------------------------------
%  Schemas TikZ du dossier T005
%  Unite : le millimetre. La largeur utile du texte vaut 160 mm.
% ---------------------------------------------------------------------------

% Ligne secondaire d'un noeud : la fonte doit etre reappliquee ligne par ligne,
% car l'alignement d'un noeud TikZ place chaque ligne dans son propre groupe.
\newcommand{\fline}[1]{{\tiny\ttfamily #1}}

% ------------------------------------------------------- figure : architecture
\newcommand{\figArchitecture}{\resizebox{\textwidth}{!}{%
\begin{tikzpicture}[x=1mm,y=1mm,every node/.style={font=\scriptsize}]

  % ---- panneau 1 : repartition physique
  \node[titrepan] at (0,93) {Panneau 1 \quad Répartition des quatre instances
                             sur les trois hôtes};

  \foreach \x/\nom/\mgmt in {0/H1/10.30.0.11, 54/H2/10.30.0.12, 108/H3/10.30.0.13}{
    \draw[hote] (\x,58) rectangle ++(50,30);
    \node[font=\scriptsize\bfseries,text=gdnavy,anchor=north west] at (\x+2,87) {Hôte~\nom};
    \node[font=\tiny\ttfamily,text=gdgrey,anchor=north east] at (\x+48,86.5) {\mgmt};
  }

  \node[blocteal,minimum width=45mm,minimum height=9mm,anchor=south west] at (2.5,71)
    {\textbf{ns1} \,\ autoritatif primaire\\[-2pt]{\tiny\ttfamily 10.30.4.10}};
  \node[blocteal,minimum width=45mm,minimum height=9mm,anchor=south west] at (2.5,60)
    {\textbf{rec2} \,\ résolveur récursif\\[-2pt]{\tiny\ttfamily 10.30.4.13}};
  \node[blocteal,minimum width=45mm,minimum height=9mm,anchor=south west] at (56.5,71)
    {\textbf{rec1} \,\ résolveur récursif\\[-2pt]{\tiny\ttfamily 10.30.4.12}};
  \node[blocteal,minimum width=45mm,minimum height=9mm,anchor=south west] at (110.5,71)
    {\textbf{ns2} \,\ autoritatif secondaire\\[-2pt]{\tiny\ttfamily 10.30.4.11}};

  \node[font=\tiny,text=gdgrey,anchor=north west,text width=158mm] at (0,56)
    {Anti-affinité par rôle : les deux membres d'une même paire ne résident jamais
     sur le même hôte, de sorte que toute perte d'un hôte laisse un serveur de
     chaque rôle en service.};

  \draw[gdrule,line width=0.4pt] (0,49) -- (158,49);

  % ---- panneau 2 : chemins logiques
  \node[titrepan] at (0,45) {Panneau 2 \quad Chemins de résolution};

  % chemin direct interdit, place dans un couloir degage au-dessus des blocs
  \node[etiqr,font=\tiny\bfseries] at (79,39.8)
    {accès direct d'un VNet tenant aux serveurs autoritatifs : refusé au pare-feu};
  \draw[draw=gdred,line width=0.65pt,dash pattern=on 2.4pt off 1.8pt]
       (20,32) -- (20,36.4) -- (140,36.4) -- (140,32);
  \draw[gdred,line width=0.9pt] (77.8,35.2) -- (80.2,37.6);
  \draw[gdred,line width=0.9pt] (77.8,37.6) -- (80.2,35.2);

  \node[bloc,minimum width=40mm,minimum height=16mm,anchor=south west] (clients) at (0,16)
    {\textbf{Clients autorisés}\\[1pt]
     \fline{VNets tenant 10.30.64.0/18}\\
     \fline{gestion 10.30.0.0/24}\\
     \fline{VPN admin 10.30.16.32/28}};
  \node[blocteal,minimum width=40mm,minimum height=16mm,anchor=south west] (rec) at (60,16)
    {\textbf{Résolveurs rec1 et rec2}\\[1pt]
     \fline{10.30.4.12 et .13}\\
     \fline{allow-from puis vue}\\
     \fline{cache indexé par étiquette}};
  \node[bloc,minimum width=40mm,minimum height=16mm,anchor=south west] (auth) at (120,16)
    {\textbf{Autoritatifs ns1 et ns2}\\[1pt]
     \fline{10.30.4.10 et .11}\\
     \fline{zone gandal.internal}\\
     \fline{et 30.10.in-addr.arpa}};

  \draw[fl] (clients) -- node[etiqt,above=1pt,font=\tiny\ttfamily]{UDP/TCP 53} (rec);
  \draw[flnavy] (rec) -- node[etiq,above=1pt,font=\tiny\ttfamily]{forward-zones} (auth);

  \node[blocgris,minimum width=40mm,minimum height=12mm,anchor=south west] (racine) at (60,0)
    {\textbf{Racine publique}\\[1pt]
     \fline{via PF-WAN puis R1}\\
     \fline{dnssec=validate}};
  \node[blocgold,minimum width=40mm,minimum height=12mm,anchor=south west] (repl) at (120,0)
    {\textbf{Réplication}\\[1pt]
     \fline{NOTIFY puis AXFR ou IXFR}\\
     \fline{signés TSIG, ns1 vers ns2}};

  \draw[flgris] (rec) -- (racine);
  \draw[flgold] (auth) -- (repl);

  \node[font=\tiny,text=gdgrey,anchor=north west,text width=158mm] at (0,-4)
    {La carte réseau unique par hôte et le commutateur SW1 unique restent des
     dépendances communes. La redondance couvre la perte d'une machine virtuelle
     ou d'un hôte, et non celle de SW1 ou du routeur R1.};
\end{tikzpicture}}}

% -------------------------------------------------------------- figure : zones
\newcommand{\figZones}{\resizebox{\textwidth}{!}{%
\begin{tikzpicture}[x=1mm,y=1mm,every node/.style={font=\scriptsize}]

  \node[titrepan] at (0,62) {Espace direct};
  \node[titrepan] at (84,62) {Espace inverse};

  \node[bloc,minimum width=74mm,minimum height=12mm,anchor=north west] (gi) at (0,58)
    {\textbf{gandal.internal}\\[1pt]
     \fline{SOA, NS ns1 et ns2, adresses des 4 serveurs}\\
     \fline{délégations NS vers les zones filles}};
  \node[blocteal,minimum width=66mm,minimum height=10mm,anchor=north west] (gia) at (8,41)
    {\textbf{infra.gandal.internal}\\[1pt]
     \fline{hôtes, XO, bastion, pfSense, services}};
  \node[blocteal,minimum width=66mm,minimum height=10mm,anchor=north west] (gib) at (8,27)
    {\textbf{tenant-001 à tenant-050}\\[1pt]
     \fline{50 zones, une par tenant}};
  \draw[gdgrey,line width=0.4pt] (3.5,46) |- (gia.west);
  \draw[gdgrey,line width=0.4pt] (3.5,46) |- (gib.west);

  \node[bloc,minimum width=74mm,minimum height=12mm,anchor=north west] (ri) at (84,58)
    {\textbf{30.10.in-addr.arpa}\\[1pt]
     \fline{le /16 appartient entièrement au projet}\\
     \fline{aucun découpage RFC 2317 nécessaire}};
  \node[blocteal,minimum width=66mm,minimum height=10mm,anchor=north west] (ria) at (92,41)
    {\textbf{11 zones d'infrastructure}\\[1pt]
     \fline{octets 0, 1, 2, 4, 5, 6, 7, 8, 10, 16 et 17}};
  \node[blocteal,minimum width=66mm,minimum height=10mm,anchor=north west] (rib) at (92,27)
    {\textbf{50 zones tenant}\\[1pt]
     \fline{octets 64 à 113, une zone par /24}};
  \draw[gdgrey,line width=0.4pt] (87.5,46) |- (ria.west);
  \draw[gdgrey,line width=0.4pt] (87.5,46) |- (rib.west);

  \node[bande,minimum width=155mm,anchor=north west,text width=150mm,align=left] at (0,12)
    {{\tiny\bfseries\scshape\color{gdteal} Correspondance entre un tenant, sa zone
      directe et sa zone inverse}\\[3pt]
     $\mathcal{D}(n)=\texttt{tenant-0}n\texttt{.gandal.internal}$ \quad et \quad
     $\mathcal{R}(n)=(63+n)\texttt{.30.10.in-addr.arpa}$, \quad
     pour $1\leqslant n\leqslant 50$\\[3pt]
     {\tiny Exemple pour $n=7$ : \texttt{tenant-007.gandal.internal} et
      \texttt{70.30.10.in-addr.arpa}, préfixe \texttt{10.30.70.0/24}.}};
\end{tikzpicture}}}

% ----------------------------------------------------------- figure : sequence
\newcommand{\figSequence}{\resizebox{\textwidth}{!}{%
\begin{tikzpicture}[x=1mm,y=1mm,every node/.style={font=\scriptsize}]

  \foreach \x/\y/\n/\ti/\la/\lb in {%
     0/44/1/{Bail attribué}/{Kea délivre 10.30.70.37}/{à la VIF de la machine},
    56/44/2/{Publication du A}/{agent vers l'API de ns1}/{PATCH rrset, type REPLACE},
   112/44/3/{Publication du PTR}/{zone 70.30.10.in-addr.arpa}/{même operation\_id},
     0/20/4/{Réplication}/{NOTIFY de ns1 vers ns2}/{AXFR ou IXFR signé TSIG},
    56/20/5/{Contrôle}/{requête sur ns1, ns2,}/{rec1 et rec2 ; concordance},
   112/20/6/{État READY}/{exposé au Backend}/{operation\_id conservé}}{%
    \draw[blocteal,fill=white] (\x,\y) rectangle ++(46,17);
    \fill[gdteal] (\x,\y+13) rectangle ++(46,4);
    \node[font=\tiny\bfseries,text=white,anchor=west]  at (\x+1.6,\y+15) {\n};
    \node[font=\tiny\bfseries,text=white]              at (\x+25,\y+15) {\ti};
    \node[font=\tiny\ttfamily,text=gdgrey]             at (\x+23,\y+8.5) {\la};
    \node[font=\tiny\ttfamily,text=gdgrey]             at (\x+23,\y+4.5) {\lb};
  }

  \draw[flnavy] (46,50) -- (56,50);
  \draw[flnavy] (102,50) -- (112,50);
  \draw[flnavy] (46,26) -- (56,26);
  \draw[flnavy] (102,26) -- (112,26);
  \draw[flnavy] (135,44) -- (135,40) -- (23,40) -- (23,37);

  \draw[flrouge] (79,20) -- (79,15);
  \draw[blocrouge,fill=gdwarn] (14,4) rectangle ++(130,11);
  \node[font=\tiny\bfseries,text=gdred] at (79,11.5)
    {Échec de publication ou de contrôle : l'opération reste observable en état dégradé,};
  \node[font=\tiny,text=gdred] at (79,7)
    {l'adresse n'est pas rendue au pool et une alerte est levée.};
\end{tikzpicture}}}

% ------------------------------------------------------------- figure : chemin
\newcommand{\figChemin}{\resizebox{\textwidth}{!}{%
\begin{tikzpicture}[x=1mm,y=1mm,every node/.style={font=\scriptsize}]

  \node[bloc,minimum width=34mm,minimum height=13mm,anchor=south west] (vm) at (0,64)
    {\textbf{VM du tenant 007}\\[1pt]
     \fline{10.30.70.37}\\
     \fline{resolv.conf : .12 puis .13}};
  \node[bloc,minimum width=34mm,minimum height=13mm,anchor=south west] (gw) at (41,64)
    {\textbf{Passerelle .1}\\[1pt]
     \fline{paire PF-USER}\\
     \fline{VIP 10.30.70.1}};
  \node[blocrouge,minimum width=34mm,minimum height=13mm,anchor=south west] (pf) at (82,64)
    {\textbf{Filtre pfSense}\\[1pt]
     \fline{53 vers .12 et .13}\\
     \fline{.10 et .11 refusés}};
  \node[blocteal,minimum width=34mm,minimum height=13mm,anchor=south west] (re) at (123,64)
    {\textbf{rec1 ou rec2}\\[1pt]
     \fline{allow-from vérifié}\\
     \fline{gettag() calcule la vue}};
  \draw[fl] (vm) -- (gw);   \draw[fl] (gw) -- (pf);   \draw[fl] (pf) -- (re);

  \node[bande,minimum width=155mm,anchor=north west,text width=150mm,align=left] at (0,58)
    {{\tiny\bfseries\scshape\color{gdteal} preresolve() : politique de sélection
      des zones}\\[3pt]
     {\tiny\ttfamily nom sous tenant-007.gandal.internal ou 70.30.10.in-addr.arpa
      \ :\ résolution autorisée}\\[1.5pt]
     {\tiny\ttfamily\color{gdred} nom sous infra.gandal.internal ou sous un autre
      tenant \ :\ REFUSED, sans divulgation}\\[1.5pt]
     {\tiny\ttfamily\color{gdgrey} nom public \ :\ récursion ordinaire, quelle que
      soit la vue}};
  \draw[fl] (140,64) -- (140,58);

  \node[bloc,minimum width=46mm,minimum height=13mm,anchor=south west] (na) at (0,14)
    {\textbf{Autoritatifs ns1 et ns2}\\[1pt]
     \fline{10.30.4.10 et .11}\\
     \fline{atteints par forward-zones}};
  \node[blocgris,minimum width=46mm,minimum height=13mm,anchor=south west] (rp) at (56,14)
    {\textbf{Racine publique}\\[1pt]
     \fline{via PF-WAN puis R1}\\
     \fline{dnssec=validate}};
  \node[blocgold,minimum width=46mm,minimum height=13mm,anchor=south west] (jo) at (112,14)
    {\textbf{Journal dnstap}\\[1pt]
     \fline{horodatage en UTC}\\
     \fline{lu en pull par Monitoring}};
  \draw[flnavy] (23,34) -- (23,27);
  \draw[flgris] (79,34) -- (79,27);
  \draw[flgold] (135,34) -- (135,27);

  \node[font=\tiny,text=gdgrey,anchor=north west,text width=158mm] at (0,12)
    {\textbf{\color{gdteal}L'étiquette de vue est calculée avant toute consultation
     de cache et entre dans la clé du cache paquet : deux vues distinctes ne peuvent
     jamais partager une réponse.} Un contrôle placé dans \texttt{preresolve()} seul
     n'offrirait pas cette garantie.};
\end{tikzpicture}}}


\begin{document}

% ---------------------------------------------------------------- couverture
\begin{titlepage}
\thispagestyle{empty}
\setlength{\parindent}{0pt}
\begin{tikzpicture}[remember picture,overlay]
  \fill[gdnavy] (current page.north west) rectangle
       ($(current page.south west)+(11mm,0)$);
  \fill[gdgold] ($(current page.north west)+(11mm,-62mm)$) rectangle
       ++(3.2mm,-46mm);
  \draw[gdrule,line width=0.8pt]
       ($(current page.north east)+(-60mm,-20mm)$) --
       ($(current page.north east)+(-16mm,-64mm)$) --
       ($(current page.north east)+(-16mm,-20mm)$) -- cycle;
  \draw[gdgold,line width=0.7pt]
       ($(current page.north east)+(-49mm,-31mm)$) --
       ($(current page.north east)+(-24mm,-56mm)$) --
       ($(current page.north east)+(-24mm,-31mm)$);
\end{tikzpicture}

\vspace*{18mm}
{\color{gdnavy}\bfseries\large Université de Yaoundé I\par}
{\color{gdnavy}\bfseries\large École Nationale Supérieure Polytechnique de Yaoundé\par}
\vspace{1.5mm}
{\color{gdgrey}\normalsize Département de Génie Informatique\par}

\vspace{40mm}

{\scshape\bfseries\color{gdteal}\small
 Dossier de réalisation GANDAL \quad\textbar\quad Tâche T005\par}
\vspace{5mm}

{\color{gdnavy}\bfseries\fontsize{30}{36}\selectfont
 Mise en place\\[2mm] du service DNS\par}

\vspace{6mm}
{\color{gdgrey}\large\itshape
 Conception détaillée, configurations et procédures\\
 d'exploitation du socle de résolution de noms\par}

\vspace{7mm}
{\color{gdgold}\rule{58mm}{1.4pt}\par}

\vspace{16mm}
{\color{gdgrey}\small Présenté dans le cadre de\par}
\vspace{1mm}
{\color{gdnavy}\large Spécifications d'ingénierie datacenter\par}

\vfill

\noindent
\begin{minipage}[t]{0.52\textwidth}
  {\scriptsize\scshape\bfseries\color{gdteal} Réalisé par\par}
  \vspace{1.5mm}
  {\color{gdnavy}\large FOMETHE SOMBANANG Maximilien\par}
  {\color{gdgrey}\small Équipe 20, projet GANDAL\par}
\end{minipage}%
\begin{minipage}[t]{0.48\textwidth}
  {\scriptsize\scshape\bfseries\color{gdteal} Périmètre\par}
  \vspace{1.5mm}
  {\color{gdnavy}\large Cellule Réseaux et Sécurité\par}
  {\color{gdgrey}\small Lot Services Network, tâche T005\par}
\end{minipage}

\vspace{7mm}
{\color{gdrule}\rule{\textwidth}{0.6pt}\par}
\vspace{2mm}
\noindent
{\color{gdgrey}\small Année universitaire 2025 / 2026 \hfill
 Version 1.0 \quad\textbar\quad 5 octobre 2026\par}
\end{titlepage}

% --------------------------------------------------------- pieces liminaires
\pagenumbering{roman}
\setcounter{page}{1}

% =========================== Fiche de tache, resume et abreviations =========
\chapter*{Fiche de tâche et résumé}
\markboth{Fiche de tâche et résumé}{}
\addcontentsline{toc}{chapter}{Fiche de tâche et résumé}

\begin{table}[H]
\centering
\small
\begin{tabular}{@{}R{0.26}R{0.74}@{}}
\hautfilet
\textbf{Rubrique} & \textbf{Contenu}\\
\medfilet
Identifiant & T005\\
Intitulé & Mettre en place le DNS\\
Responsable & FOMETHE SOBMBANANG MAXIMILIEN\\
Pondération et priorité & 8 sur 10, priorité haute\\
Cellule & Réseaux et Sécurité, lot Services Network\\
Prérequis consommés & T001, plan d'adressage et modèle IPAM figés par le
  document 02 ; T003, réseau de management\\
Tâches coordonnées & T004 DHCP haute disponibilité, T006 synchronisation
  temporelle, T007 agent réseau et API de provisionnement\\
Livrable & Le présent dossier : conception, configurations, procédures et
  critères de recette du service DNS\\
Hors périmètre & T009, tests, validation, reprise et documentation, qui
  exploite les critères définis au chapitre~\ref{ch:recette} une fois les
  autres lots disponibles\\
\basfilet
\end{tabular}
\end{table}

\section*{Résumé}

Ce dossier définit et paramètre le service de résolution de noms de GANDAL. Il
sépare strictement le rôle autoritatif du rôle récursif, retient PowerDNS
conformément à l'arbitrage du document 03, et construit l'espace de nommage
interne à partir du plan d'adressage arrêté par le document 02. Il traite la
confidentialité des noms entre domaines par une politique de sélection des
zones indexée sur l'adresse source, doublée d'un filtrage réseau qui interdit
aux VNets tenant d'atteindre les serveurs autoritatifs.

Le cycle de vie des enregistrements est confié à un propriétaire d'écriture
unique, l'agent réseau, au moyen de l'API autoritative exposée derrière une
terminaison TLS. Les opérations sont rejouables sans effet de bord, et une
panne de publication laisse l'opération dans un état dégradé observable sans
libérer une adresse encore utilisée.

Les configurations proposées sont complètes et cohérentes avec les contraintes
matérielles du projet : une carte réseau par hôte, un commutateur unique, une
MTU utile de 1\,450 octets imposée par l'encapsulation VXLAN. Conformément
à la règle du corpus, aucun déploiement ni aucun essai n'est présenté comme
effectué : les contrôles de recette sont énoncés avec leur commande, leur
résultat attendu et le statut \statut{non exécuté}.

\begin{encadre}
\titreencadre{Mots-clés}
GANDAL ; DNS ; PowerDNS ; zone autoritative ; résolveur récursif ; espace de
nommage interne ; isolation multi-tenant ; TSIG ; DNSSEC ; publication
idempotente ; mode dégradé.
\end{encadre}

\section*{Abréviations propres au dossier}

Le glossaire général du corpus reste applicable. Les termes ci-dessous sont
ceux que ce dossier emploie de manière spécifique.

\begin{table}[H]
\centering
\small
\begin{tabular}{@{}R{0.17}R{0.83}@{}}
\hautfilet
\textbf{Terme} & \textbf{Définition retenue}\\
\medfilet
Autoritatif & Serveur qui détient les données d'une zone et répond avec
  l'indicateur \opt{AA}. Il ne résout pas les noms des autres zones.\\
Récursif & Résolveur qui interroge la chaîne des serveurs autoritatifs pour le
  compte d'un client et met les réponses en cache.\\
AXFR, IXFR & Transfert complet ou incrémental d'une zone entre serveurs.\\
TSIG & Signature symétrique des messages DNS, définie par le RFC 8945,
  employée ici pour les transferts de zone.\\
NTA & \textit{Negative Trust Anchor} : exception locale de validation DNSSEC
  sur une branche de l'arbre.\\
RRset & Ensemble des enregistrements partageant le même nom, la même classe et
  le même type ; unité de modification de l'API.\\
PTR & Enregistrement de résolution inverse, d'une adresse vers un nom.\\
Vue & Ensemble de zones qu'une catégorie de clients est autorisée à résoudre ;
  déterminée ici par l'adresse source.\\
dnstap & Format et transport de journalisation structurée des échanges DNS.\\
EDNS & Extension du protocole DNS permettant notamment d'annoncer une taille de
  réponse UDP supérieure à 512 octets.\\
\basfilet
\end{tabular}
\end{table}

\section*{Notations}

Le dossier emploie les notations suivantes, dont les définitions complètes
figurent aux sections~\ref{sec:nommage} et~\ref{sec:dimensionnement}.

\begin{table}[H]
\centering
\small
\begin{tabular}{@{}R{0.13}R{0.87}@{}}
\hautfilet
\textbf{Symbole} & \textbf{Signification}\\
\medfilet
$n$ & Numéro d'un tenant, $1\leqslant n\leqslant T$\\
$T$ & Nombre de tenants du périmètre minimal, $T=50$\\
$\mathcal{P}(n)$ & Préfixe IPv4 attribué au tenant $n$\\
$\mathcal{D}(n)$, $\mathcal{R}(n)$ & Zone directe et zone inverse du tenant $n$\\
$\tau(a)$ & Étiquette de vue associée à l'adresse source $a$\\
$o_3(a)$ & Troisième octet de l'adresse $a$\\
$N_Z$, $N_R$ & Nombre de zones et nombre d'enregistrements publiés\\
$B$, $Q$ & Durée du bail DHCP et durée de quarantaine d'une adresse\\
\basfilet
\end{tabular}
\end{table}


\cleardoublepage
\phantomsection
\addcontentsline{toc}{chapter}{Table des matières}
\renewcommand{\contentsname}{Table des matières}
{\hypersetup{linkcolor=gdnavy}\tableofcontents}

\cleardoublepage
\pagenumbering{arabic}
\setcounter{page}{1}

% ------------------------------------------------------------------ chapitres
% ================== Chapitre 1 : objectifs, perimetre et exigences ==========
\chapter{Objectifs, périmètre et exigences couvertes}
\label{ch:objectifs}

\section{Objectif de la tâche et définition du service attendu}

La tâche T005 consiste à mettre en place le service DNS de GANDAL. Dans le
découpage du suivi d'équipe, elle se situe après le gel du plan d'adressage et
du modèle IPAM, et en parallèle du DHCP haute disponibilité et de la
synchronisation temporelle. Le service attendu n'est pas un simple serveur de
noms : le document 03 lui assigne quatre fonctions distinctes.

\begin{itemize}
  \item \textbf{Publier l'espace de nommage interne du datacenter}, c'est-à-dire
    les noms d'infrastructure et les noms des machines de chaque tenant, avec
    les résolutions inverses correspondantes.
  \item \textbf{Fournir une résolution récursive contrôlée} aux machines
    autorisées, pour les noms internes comme pour les noms publics.
  \item \textbf{Garantir la confidentialité des noms entre domaines} : un tenant
    ne doit lire ni la zone d'infrastructure ni la zone d'un autre tenant.
  \item \textbf{S'insérer dans le cycle de vie d'une allocation} : la publication
    d'un enregistrement est une étape de la machine d'états de création, entre
    l'attribution du bail et l'exposition de l'état \opt{READY}.
\end{itemize}

Mettre en place le DNS signifie donc arrêter une architecture, un espace de
nommage, des règles d'accès et un mécanisme de mise à jour, puis produire les
configurations et procédures qui permettent de les appliquer de manière
reproductible. C'est l'objet des chapitres \ref{ch:conception} à
\ref{ch:cycle}.

\begin{encadre}
\titreencadre{Règle de restitution du corpus}
Tous les documents de cadrage de GANDAL précisent que la règle est d'effectuer
une simulation au préalable et qu'aucun déploiement ni essai n'est présenté
comme effectué. Ce dossier respecte cette règle : il fournit des configurations
et des procédures complètes, mais les contrôles du chapitre~\ref{ch:recette}
portent le statut \statut{non exécuté}. Produire les preuves correspondantes
relève de la tâche T009, qui n'est pas traitée ici.
\end{encadre}

\section{Contexte repris du corpus et données d'entrée}

Les choix qui suivent sont contraints par des décisions déjà arrêtées. Elles ne
sont pas rediscutées ici ; elles sont rappelées parce qu'elles déterminent
directement la conception du service.

\begin{longtable}{@{}R{0.31}R{0.52}R{0.17}@{}}
\caption{Données d'entrée de la tâche T005 et leurs conséquences directes}
\label{tab:entrees}\\
\hautfilet
\textbf{Contrainte ou décision antérieure} &
\textbf{Conséquence sur le service DNS} & \textbf{Source}\\
\medfilet
\endfirsthead
\hautfilet
\textbf{Contrainte ou décision antérieure} &
\textbf{Conséquence sur le service DNS} & \textbf{Source}\\
\medfilet
\endhead
\basfilet
\endfoot
Trois hôtes XCP-ng, une NIC par hôte, un commutateur L2 et un routeur non
  manageables &
  La redondance logicielle ne peut couvrir que la perte d'une machine virtuelle
  ou d'un hôte. SW1 et R1 restent des dépendances communes assumées. &
  Document 01, 2.1\\
Bloc \ip{10.30.0.0/16}, VNet Services en \ip{10.30.4.0/24}, DNS autoritatif en
  \ip{.10} et \ip{.11}, DNS récursif en \ip{.12} et \ip{.13} &
  Les adresses des quatre instances sont imposées. Ce dossier les consomme sans
  redéfinition locale. & Document 02, tableau 2.5\\
Tenant $n$ associé au préfixe \ip{10.30.(63+n).0/24}, pour
  $1\leqslant n\leqslant 50$ &
  La correspondance entre un tenant, sa zone directe et sa zone inverse se
  déduit arithmétiquement, ce qui rend la génération des zones et la politique
  de vue déterministes. & Document 02, 2.4\\
VXLAN sur IPv4 avec MTU externe de 1\,500 octets &
  Au plus 1\,450 octets utiles pour le paquet IP interne. Le dimensionnement
  des réponses DNS en doit tenir compte. & Document 02, 2.5\\
PowerDNS Authoritative complété de PowerDNS Recursor &
  Choix proposé par la comparaison du document 03 ; BIND reste l'alternative
  documentée du corpus. & Document 03, tableau 1.1\\
Aucun VNet surveillé n'initie de connexion vers Monitoring &
  Les journaux et métriques du DNS sont conservés localement et lus en pull.
  Aucun envoi direct n'est configuré. & Document 03, 2.7\\
Deux instances de chaque service, sur des hôtes distincts &
  Objectif de continuité logique, sans prétendre à l'élimination des points
  uniques de défaillance. & \exig{net-nfr-001}\\
\end{longtable}

\section{Registre des exigences couvertes}

Le tableau \ref{tab:tracabilite} relie chaque exigence du registre des services
réseau au mécanisme retenu dans ce dossier et au contrôle de recette qui
permettra de la vérifier. Les identifiants de contrôle \opt{C01} à \opt{C15}
sont définis au chapitre~\ref{ch:recette}.

\begin{longtable}{@{}R{0.16}R{0.18}R{0.50}R{0.16}@{}}
\caption{Traçabilité entre les exigences du registre et les mécanismes de ce
  dossier}\label{tab:tracabilite}\\
\hautfilet
\textbf{Exigence} & \textbf{Objet} & \textbf{Mécanisme retenu et section} &
\textbf{Contrôle}\\
\medfilet
\endfirsthead
\hautfilet
\textbf{Exigence} & \textbf{Objet} & \textbf{Mécanisme retenu et section} &
\textbf{Contrôle}\\
\medfilet
\endhead
\basfilet
\endfoot
NET-FR-004 & Paramètres DNS par tenant & Objet profil DNS porté par l'IPAM et
  consommé par DHCP et Templates (\ref{sec:profil}) & C03\\
NET-FR-006 & DNS interne et mises à jour & Zones internes autoritatives et
  publication par l'agent via l'API (\ref{sec:nommage}, \ref{sec:publication})
  & C01, C10\\
NET-FR-007 & Récursion DNS autorisée & \opt{allow-from} restreint aux préfixes
  du plan, sortie 53 explicite (\ref{sec:recursif}) & C07\\
NET-FR-008 & Confidentialité des noms entre domaines & Politique de vue par
  adresse source, transferts fermés, autoritatif non exposé
  (\ref{sec:vue}, \ref{sec:flux}) & C04, C05, C06\\
NET-FR-012 & Allocation et libération idempotentes & Modification par RRset avec
  lecture préalable et clé d'idempotence (\ref{sec:idempotence}) & C10\\
NET-FR-013 & Paramètres VM via Templates & Profil DNS transmis par le contrat
  IF-NS-06 (\ref{sec:profil}) & C03\\
NET-FR-014 & Métriques services & Point de métriques en lecture seule,
  interrogé en pull (\ref{sec:journaux}) & C14\\
NET-FR-016 & API administratives chiffrées & API autoritative liée à la boucle
  locale derrière une terminaison TLS (\ref{sec:api}) & C10\\
NET-FR-017 & DNS accessible selon profil VPN & Vue administration pour le pool
  VPN admin, refus par défaut pour le pool utilisateur (\ref{sec:vpn}) & C06\\
NET-FR-018 & Journaux horodatés & Journalisation dnstap horodatée en temps
  universel, conservée localement (\ref{sec:journaux}) & C14\\
NET-FR-019 & Gestion de l'infrastructure isolée & Les VNets tenant ne joignent
  ni le plan de gestion ni les autoritatifs (\ref{sec:flux}) & C08\\
NET-FR-020 & Simulation avant déploiement & Statut \statut{non exécuté}
  maintenu ; maquette préalable exigée (chapitre~\ref{ch:limites}) & aucun\\
NET-NFR-001 & Deux instances DNS & Deux autoritatifs et deux récursifs répartis
  sur les trois hôtes (\ref{sec:placement}) & C09\\
NET-NFR-003 & Reconstruction en moins de 60 min & Export des zones et procédure
  de restauration (\ref{sec:sauvegarde}) ; mesure confiée à T009 & aucun\\
NET-NFR-005 & Protection des communications administratives & TLS sur l'API,
  TSIG sur les transferts (\ref{sec:zones}, \ref{sec:api}) & C04\\
NET-NFR-006 & Rejeu déterministe et idempotent & Opérations convergentes sur
  l'état désiré des RRsets (\ref{sec:idempotence}) & C10\\
NET-NFR-007 & Secrets absents des dépôts en clair & Clés API et TSIG injectées
  au démarrage ; empreinte en configuration (\ref{sec:sauvegarde}) & aucun\\
IF-NS-03 & DNS vers VNet et SDN : profils et ACL par tenant & Étiquette de vue
  dérivée du préfixe tenant (\ref{sec:vue}) & C05\\
IF-NS-06 & Agent vers Templates : réseau, DNS, MTU & Contenu normalisé du
  profil DNS (\ref{sec:profil}) & C03\\
IF-NS-08 & Network vers Monitoring : lecture en pull & Aucune initiation
  sortante depuis les VM DNS (\ref{sec:journaux}) & C14\\
\end{longtable}

\section{Limites de périmètre et articulation avec les autres tâches}

Trois tâches voisines touchent au même cycle de vie sans relever de T005. Les
points de contact sont explicités pour éviter une double définition.

\begin{longtable}{@{}R{0.30}R{0.70}@{}}
\caption{Frontières de la tâche T005}\label{tab:frontieres}\\
\hautfilet
\textbf{Tâche voisine} & \textbf{Point de contact et répartition}\\
\medfilet
\endfirsthead
\hautfilet
\textbf{Tâche voisine} & \textbf{Point de contact et répartition}\\
\medfilet
\endhead
\basfilet
\endfoot
T004, DHCP haute disponibilité \newline (B.~Heudep Djandja) &
  T005 fournit les adresses des résolveurs et le nom de domaine à distribuer
  dans les options 6, 15 et 119. T005 ne définit ni les scopes, ni le mode de
  haute disponibilité de Kea, ni la durée de bail. La valeur de bail est ici
  une hypothèse de travail, signalée comme telle.\\
T006, synchronisation temporelle \newline (B.~Heudep Djandja) &
  La validation DNSSEC et la journalisation corrélable supposent une horloge
  correcte. T005 énonce cette dépendance et la contrainte d'amorçage à froid,
  mais ne configure pas chrony.\\
T007, agent réseau et API de provisionnement \newline (P.~Pangui Pianne) &
  T005 définit le contrat que l'agent doit respecter pour publier un
  enregistrement et le comportement attendu en cas d'échec. L'implémentation de
  l'agent ne relève pas de ce dossier.\\
T009, tests, validation, reprise et documentation &
  Seconde tâche attribuée au même responsable. Elle dépend des résultats des
  lots précédents et n'est pas traitée ici. Le chapitre~\ref{ch:recette} lui
  transmet des critères d'acceptation déjà formulés, qui devront être exécutés
  et archivés dans ce cadre.\\
\end{longtable}

% ================== Chapitre 2 : conception du service de resolution ========
\chapter{Conception du service de résolution}
\label{ch:conception}

\section{Séparation des rôles autoritatif et récursif}
\label{sec:roles}

Le document 03 impose de distinguer au moins deux rôles de service et précise
qu'une instance récursive partagée non filtrée est insuffisante. Cette
séparation n'est pas une préférence de mise en \oe uvre : elle découle de trois
propriétés incompatibles sur une même instance.

\begin{itemize}
  \item Un serveur autoritatif doit répondre avec autorité pour les zones qu'il
    détient et ne doit rien résoudre d'autre. Un résolveur doit au contraire
    parcourir l'arbre public et conserver un cache.
  \item Les populations de clients diffèrent. L'autoritatif n'a que quatre
    interlocuteurs légitimes : son homologue, les deux résolveurs et l'agent. Le
    résolveur sert au contraire l'ensemble des machines autorisées.
  \item Les surfaces d'attaque diffèrent. Exposer l'autoritatif aux tenants
    reviendrait à leur offrir un point d'interrogation direct de toutes les
    zones hébergées, que seule une politique interne pourrait ensuite
    restreindre.
\end{itemize}

La conception retient donc deux serveurs autoritatifs, \opt{ns1} et \opt{ns2},
qui ne sont joignables que depuis le VNet Services, et deux résolveurs
récursifs, \opt{rec1} et \opt{rec2}, qui constituent le seul point de résolution
offert aux machines clientes. Les résolveurs atteignent l'espace interne par une
redirection explicite de zone, et l'espace public par une récursion ordinaire.

\begin{figure}[H]
\centering
\figArchitecture
\caption{Architecture du service DNS : placement des quatre instances, flux
  autorisés et dépendances communes}
\label{fig:architecture}
\end{figure}

\section{Espace de nommage, zones directes et zones inverses}
\label{sec:nommage}

La zone \zone{gandal.internal} est conservée comme proposition de nommage
interne, conformément au document 03. Le suffixe est pertinent : l'ICANN a
réservé le domaine de premier niveau \zone{internal} à l'usage privé et ne le
délègue pas dans la racine publique. Un nom construit sous ce suffixe ne peut
donc pas entrer en collision avec un nom Internet, ce qui n'est pas le cas d'un
suffixe arbitraire.

\subsection{Découpage retenu}

\begin{itemize}
  \item \zone{gandal.internal} porte le SOA, les enregistrements NS du service
    et les adresses des quatre serveurs, puis délègue ses zones filles.
  \item \zone{infra.gandal.internal} contient l'inventaire d'infrastructure :
    hôtes, orchestrateur, bastion, machines pfSense, services de base et
    collecteurs.
  \item \zone{tenant-0NN.gandal.internal}, pour \opt{NN} de \opt{001} à
    \opt{050}, contient la passerelle du tenant et les noms de ses machines.
\end{itemize}

Les serveurs de noms sont volontairement nommés au niveau de l'apex, sous la
forme \zone{ns1.gandal.internal}, et non à l'intérieur de
\zone{infra.gandal.internal}. Ce choix évite une délégation dont les serveurs
porteraient un nom situé sous la zone déléguée, configuration qui exige des
enregistrements de collage dans la zone parente et complique inutilement la
reconstruction.

Le plan d'adressage du document 02 rend la correspondance entre un tenant et ses
zones entièrement déterministe. En notant $T=50$ le nombre de tenants du
périmètre minimal, elle s'écrit
\begin{align}
  \mathcal{P}(n) &= \texttt{10.30.}(63+n)\texttt{.0/24},
    & 1 &\leqslant n \leqslant T, \label{eq:prefixe}\\[2pt]
  \mathcal{D}(n) &= \texttt{tenant-0}\langle n\rangle\texttt{.gandal.internal},
    \label{eq:directe}\\[2pt]
  \mathcal{R}(n) &= (63+n)\texttt{.30.10.in-addr.arpa},
    \label{eq:inverse}
\end{align}
où $\langle n\rangle$ désigne le numéro du tenant écrit sur trois chiffres. Le
tenant 1 occupe ainsi \ip{10.30.64.0/24} et le tenant 50 occupe
\ip{10.30.113.0/24}, les quatorze préfixes \ip{10.30.114.0/24} à
\ip{10.30.127.0/24} restant en réserve à l'intérieur du \ip{/18}.

\subsection{Résolution inverse}

Toutes les plages du projet appartiennent au bloc \ip{10.30.0.0/16} et sont
découpées en \ip{/24}, en \ip{/28} ou en \ip{/29}. Comme le projet détient
l'intégralité de l'espace, une zone inverse par \ip{/24} suffit à couvrir les
sous-réseaux plus petits qu'elle contient : les quatre \ip{/28} de la plage DMZ
et VPN relèvent tous de \zone{16.30.10.in-addr.arpa}, et les trois \ip{/29} de
synchronisation de \zone{17.30.10.in-addr.arpa}. Le découpage sans classe du
RFC 2317, utile lorsqu'un \ip{/24} est partagé entre plusieurs autorités, n'a
donc aucune raison d'être introduit ici.

\begin{figure}[H]
\centering
\figZones
\caption{Arborescence des zones directes et inverses}
\label{fig:zones}
\end{figure}

Le nombre total de zones à créer se déduit de ce découpage. En comptant la zone
racine interne et la zone d'infrastructure, les onze zones inverses
d'infrastructure, puis une zone directe et une zone inverse par tenant,
\begin{equation}
  N_Z = \underbrace{2}_{\text{apex et infra}}
      + \underbrace{11}_{\text{inverses d'infrastructure}}
      + \underbrace{2T}_{\text{zones des tenants}}
      = 2T + 13 = 113 .
  \label{eq:nzones}
\end{equation}

\begin{longtable}{@{}R{0.21}R{0.30}C{0.13}R{0.36}@{}}
\caption{Inventaire des zones à créer}\label{tab:zones}\\
\hautfilet
\textbf{Catégorie} & \textbf{Nom ou intervalle} & \textbf{Nombre} &
\textbf{Contenu}\\
\medfilet
\endfirsthead
\hautfilet
\textbf{Catégorie} & \textbf{Nom ou intervalle} & \textbf{Nombre} &
\textbf{Contenu}\\
\medfilet
\endhead
\basfilet
\endfoot
Zone racine interne & \zone{gandal.internal} & 1 &
  SOA, NS, adresses des serveurs, délégations\\
Zone d'infrastructure & \zone{infra.gandal.internal} & 1 &
  environ 40 noms, statiques\\
Zones tenant & \zone{tenant-001} à \zone{tenant-050} \zone{.gandal.internal} &
  50 & passerelle statique, machines publiées par l'agent\\
Inverses d'infrastructure & \zone{0, 1, 2, 4, 5, 6, 7, 8, 10, 16, 17
  .30.10.in-addr.arpa} & 11 &
  un \ip{/24} inverse couvre les \ip{/28} et \ip{/29} qu'il contient\\
Inverses tenant & \zone{64} à \zone{113} \zone{.30.10.in-addr.arpa} & 50 &
  tenant $n$ associé à l'octet $63+n$\\
\medfilet
\textbf{Total} & & \textbf{113} &
  volume compatible avec une base locale\\
\end{longtable}

\section{Placement, redondance et domaines de panne}
\label{sec:placement}

L'exigence \exig{net-nfr-001} demande deux instances par service, et le document
Underlay rappelle que la présence de plusieurs serveurs ne garantit pas leur
indépendance si tous partagent le même hyperviseur. Le placement doit donc
satisfaire une contrainte d'anti-affinité par rôle : les deux membres d'une
paire ne résident jamais sur le même hôte.

Avec quatre instances et trois hôtes, un hôte en porte nécessairement deux. La
répartition retenue place \opt{ns1} et \opt{rec2} sur H1, \opt{rec1} sur H2 et
\opt{ns2} sur H3. Elle est préférable à un regroupement sur deux hôtes
seulement, car elle laisse trois instances sur quatre en service après la perte
de H2 ou de H3.

\begin{longtable}{@{}R{0.17}R{0.17}R{0.17}R{0.49}@{}}
\caption{Effet de la perte d'un hôte sur le service de résolution}
\label{tab:pannes}\\
\hautfilet
\textbf{Scénario} & \textbf{Autoritatifs restants} &
\textbf{Récursifs restants} & \textbf{Conséquence}\\
\medfilet
\endfirsthead
\hautfilet
\textbf{Scénario} & \textbf{Autoritatifs restants} &
\textbf{Récursifs restants} & \textbf{Conséquence}\\
\medfilet
\endhead
\basfilet
\endfoot
Fonctionnement nominal & ns1 (H1), ns2 (H3) & rec1 (H2), rec2 (H1) &
  Service complet\\
Perte de H1 & ns2 & rec1 &
  Un serveur de chaque rôle, capacité réduite de moitié ; ns2 devient le seul
  autoritatif et ne reçoit plus de transfert, donc plus de publication\\
Perte de H2 & ns1, ns2 & rec2 & Résolution et publication nominales\\
Perte de H3 & ns1 & rec1, rec2 &
  Publication nominale, réplication suspendue jusqu'au retour de ns2\\
Perte de SW1 & aucune & aucune &
  Hors de portée de la redondance logicielle ; limite matérielle assumée par le
  cadrage\\
\end{longtable}

\begin{encadredecision}
\titreencadredec{Conséquence à retenir pour la reprise}
La perte de H1 supprime le seul serveur autoritatif primaire. La résolution se
poursuit sur \opt{ns2}, qui continue de servir les zones transférées jusqu'à
expiration du paramètre \opt{expire} du SOA, fixé à sept jours. En revanche,
aucune nouvelle publication n'est possible pendant cette période. La procédure
de reprise consiste à restaurer \opt{ns1} depuis l'export des zones, et non à
promouvoir \opt{ns2}, afin de conserver un propriétaire d'écriture unique.

\smallskip
Cette distinction entre continuité du trafic existant et impossibilité
d'opérations nouvelles correspond exactement au résultat attendu du test
\opt{G04} du document 01.
\end{encadredecision}

\section{Politique de durée de vie et cohérence avec le bail DHCP}
\label{sec:ttl}

Les durées de vie ne sont pas un réglage de confort : elles déterminent le délai
pendant lequel une réponse périmée peut encore circuler après la suppression
d'un enregistrement. Deux contraintes les gouvernent. La durée de vie d'un
enregistrement dynamique doit rester nettement inférieure à la durée $B$ du bail
qui l'a produit, et une adresse ne peut être réattribuée qu'après expiration des
caches qui ont pu mémoriser son ancien nom. En notant $Q$ la durée de
quarantaine et $\delta$ une marge de sécurité,
\begin{equation}
  \mathrm{TTL}_A \;\leqslant\; \frac{B}{4}
  \qquad\text{et}\qquad
  Q \;\geqslant\; \mathrm{TTL}_A + \mathrm{TTL}_{\mathrm{nég}} + \delta .
  \label{eq:ttl}
\end{equation}
Avec l'hypothèse $B = 3\,600$ secondes retenue plus bas, le choix
$\mathrm{TTL}_A = \mathrm{TTL}_{\mathrm{nég}} = 300$ secondes et
$\delta = 300$ secondes conduit à $Q = 900$ secondes et respecte les deux
inégalités avec un facteur de sécurité confortable.

\begin{longtable}{@{}R{0.28}R{0.14}R{0.58}@{}}
\caption{Politique de durée de vie}\label{tab:ttl}\\
\hautfilet
\textbf{Paramètre} & \textbf{Valeur} & \textbf{Justification}\\
\medfilet
\endfirsthead
\hautfilet
\textbf{Paramètre} & \textbf{Valeur} & \textbf{Justification}\\
\medfilet
\endhead
\basfilet
\endfoot
SOA \opt{refresh} & 3\,600 s &
  Fréquence de contrôle de ns2 en l'absence de NOTIFY ; la réplication normale
  reste déclenchée par NOTIFY.\\
SOA \opt{retry} & 900 s & Nouvelle tentative après échec de transfert.\\
SOA \opt{expire} & 604\,800 s &
  Durée pendant laquelle ns2 continue de servir une zone sans contact avec ns1 ;
  dimensionnée pour couvrir une panne longue de H1.\\
SOA \opt{minimum} & 300 s &
  Durée de mémorisation d'une réponse négative. Une valeur courte évite qu'un
  nom publié tardivement reste introuvable.\\
Enregistrements d'infrastructure & 3\,600 s &
  Noms stables, modifiés par une opération planifiée.\\
Enregistrements A et PTR tenant & 300 s &
  Publiés et retirés avec le bail, soit $B/12$ sous l'hypothèse
  $B = 3\,600$ s.\\
Quarantaine $Q$ avant réutilisation & 900 s &
  Proposée à l'IPAM. Elle satisfait l'inégalité \eqref{eq:ttl} et répond au
  point laissé \statut{à valider} par le document 03.\\
\end{longtable}

La durée de bail de 3\,600 secondes est une hypothèse de travail : elle relève
de la tâche T004. Si elle est modifiée, les deux inégalités \eqref{eq:ttl}
restent la règle à conserver.

\section{Dimensionnement et contraintes de transport}
\label{sec:dimensionnement}

\subsection{Volume d'enregistrements}

Le volume à servir se déduit du plan d'adressage. Chaque tenant dispose d'un
pool dynamique de 190 adresses, de \ip{.10} à \ip{.199}. En notant $L$ le nombre
de baux simultanés, $T$ le nombre de tenants et $H$ le nombre de noms
d'infrastructure, chaque bail produit un enregistrement A et son PTR, chaque
tenant apporte son SOA, ses deux NS et le couple A et PTR de sa passerelle, et
chaque nom d'infrastructure apporte un A et un PTR :
\begin{equation}
  N_R = 2L + 5T + 2H .
  \label{eq:nrecords}
\end{equation}
La cible historique de 500 baux simultanés, avec $T = 50$ et $H \approx 40$,
conduit à un régime courant
\[
  N_R = 2\times 500 + 5\times 50 + 2\times 40 = 1\,330
\]
enregistrements. Dans l'hypothèse haute où les cinquante pools seraient pleins,
$L_{\max} = 190\,T = 9\,500$ et
\[
  N_R^{\max} = 2\times 9\,500 + 250 + 80 = 19\,330 ,
\]
soit un peu moins de vingt mille enregistrements.

Ce volume ne justifie pas un serveur de base de données dédié. Le moteur de
stockage SQLite intégré à PowerDNS est retenu pour le périmètre minimal, la
réplication étant assurée par transfert de zone et non par duplication de base.
Un moteur relationnel partagé constitue une extension à envisager si la mesure
du débit d'écriture de l'agent le justifie.

\begin{table}[H]
\centering
\small
\caption{Dimensionnement proposé des machines virtuelles}\label{tab:vm}
\begin{tabular}{@{}R{0.14}C{0.13}C{0.12}C{0.12}R{0.49}@{}}
\hautfilet
\textbf{Instance} & \textbf{Processeur} & \textbf{Mémoire} & \textbf{Disque} &
\textbf{Remarque}\\
\medfilet
ns1, ns2 & 2 vCPU & 2 GiB & 20 GiB &
  Base SQLite, journal local, export des zones\\
rec1, rec2 & 2 vCPU & 4 GiB & 20 GiB &
  Mémoire dimensionnée pour le cache d'enregistrements et le cache paquet\\
\basfilet
\end{tabular}
\end{table}

\subsection{Taille des réponses et MTU}

Le document 02 établit la taille utile laissée au paquet IP interne par
l'encapsulation VXLAN sur IPv4, à partir d'une MTU externe de 1\,500 octets :
\begin{equation}
  \mathrm{MTU}_{\text{utile}}
    = 1\,500 - \underbrace{20}_{\text{IPv4}}
             - \underbrace{8}_{\text{UDP}}
             - \underbrace{8}_{\text{VXLAN}}
             - \underbrace{14}_{\text{Ethernet interne}}
    = 1\,450 \ \text{octets}.
  \label{eq:mtu}
\end{equation}
Une réponse DNS qui dépasserait cette taille serait fragmentée ou perdue selon
le comportement des équipements intermédiaires, et ce type de défaut se
manifeste de manière intermittente, donc difficile à diagnostiquer.

La taille annoncée par EDNS est en conséquence fixée à 1\,232 octets sur les
quatre instances. Cette valeur correspond à la recommandation issue du DNS Flag
Day 2020 et laisse une marge
\[
  \mathrm{MTU}_{\text{utile}} - 1\,232 = 218 \ \text{octets}
\]
pour les en-têtes IP et UDP du paquet interne lui-même. Au-delà de ce seuil, la
réponse est tronquée et le client bascule en TCP, comportement prévisible et
observable. Les règles de pare-feu doivent donc autoriser le port 53 en TCP
comme en UDP, point souvent omis.

\section{Profil DNS par tenant et interfaces consommatrices}
\label{sec:profil}

Le modèle IPAM minimal du document 03 comporte un champ \opt{DNS\_profile} dont
le contenu n'est pas spécifié. Ce dossier le définit, car il constitue le point
de passage unique entre la conception DNS, les options distribuées par DHCP et
les paramètres remis aux Templates.

\begin{lstlisting}[caption={Objet profil DNS porté par l'IPAM, instancié pour le
  tenant 7, de préfixe 10.30.70.0/24},label=lst:profil]
{
  "profile_id":         "tenant-standard-v1",
  "tenant_id":          7,
  "resolvers":          ["10.30.4.12", "10.30.4.13"],
  "search_domains":     ["tenant-007.gandal.internal"],
  "forward_zone":       "tenant-007.gandal.internal",
  "reverse_zone":       "70.30.10.in-addr.arpa",
  "record_ttl":         300,
  "publication":        "agent",
  "external_recursion": true,
  "view_tag":           1007
}
\end{lstlisting}

\begin{table}[H]
\centering
\small
\caption{Consommation du profil DNS par les interfaces du lot Services Network}
\label{tab:profil}
\begin{tabular}{@{}R{0.26}R{0.74}@{}}
\hautfilet
\textbf{Interface} & \textbf{Usage du profil}\\
\medfilet
IF-NS-02, vers DHCP &
  \opt{resolvers} alimente l'option 6, \opt{search\_domains} l'option 15 et
  l'option 119. Le serveur Kea ne redéfinit aucune de ces valeurs localement.\\
IF-NS-06, vers Templates &
  \opt{resolvers}, \opt{search\_domains} et la MTU utile sont remis au mécanisme
  d'initialisation de la machine. Le profil ne contient aucun secret.\\
IF-NS-03, vers VNet et SDN &
  \opt{view\_tag} matérialise le profil de résolution et la liste de contrôle
  d'accès associée au tenant.\\
IF-NS-05, vers Backend par l'agent &
  \opt{forward\_zone} et \opt{reverse\_zone} désignent les deux zones à modifier
  lors d'une publication ou d'un retrait.\\
\basfilet
\end{tabular}
\end{table}

Un second profil, \opt{admin-v1}, désigne les mêmes résolveurs avec le domaine
de recherche \zone{infra.gandal.internal} et l'étiquette de vue 1. Il s'applique
aux machines du plan de gestion et au pool VPN administrateur.

% ================== Chapitre 3 : isolation et confidentialite ===============
\chapter{Isolation et confidentialité des noms}
\label{ch:isolation}

\section{Expression du besoin et scénarios à empêcher}

L'exigence \exig{net-fr-008} porte sur la confidentialité des noms entre
domaines. Le document 03 la précise sans ambiguïté : le mécanisme de résolution
doit empêcher la lecture d'une zone d'infrastructure ou d'un autre tenant par
les clients non autorisés, et cacher un enregistrement dans le portail ne suffit
pas. Le test \opt{S04} formule la vérification attendue : depuis une machine du
VNet USER, une requête vers la zone d'infrastructure et vers la zone d'un tenant
voisin doit se solder par un refus ou une absence d'information.

Trois scénarios doivent être couverts, et chacun appelle une contre-mesure
distincte.

\begin{table}[H]
\centering
\small
\caption{Scénarios de divulgation et contre-mesures}\label{tab:scenarios}
\begin{tabular}{@{}R{0.30}R{0.33}R{0.37}@{}}
\hautfilet
\textbf{Scénario} & \textbf{Ce qui serait divulgué} &
\textbf{Contre-mesure retenue}\\
\medfilet
Interrogation directe d'un serveur autoritatif par une machine tenant &
  Le serveur autoritatif répondrait pour toute zone hébergée, y compris celles
  des autres tenants. &
  Filtrage réseau : les adresses \ip{10.30.4.10} et \ip{10.30.4.11} ne sont pas
  joignables depuis un VNet tenant (section~\ref{sec:flux}).\\
Transfert de zone demandé par un client quelconque &
  Un transfert divulgue l'intégralité d'une zone en une requête. &
  Transferts restreints à ns2 par adresse et signés par une clé TSIG
  (section~\ref{sec:zones}).\\
Requête vers une zone étrangère au travers du résolveur partagé &
  Le résolveur résout toutes les zones internes pour tous ses clients. &
  Politique de sélection des zones indexée sur l'adresse source
  (section~\ref{sec:vue}).\\
\basfilet
\end{tabular}
\end{table}

Une énumération par requêtes successives reste théoriquement possible à
l'intérieur de la vue d'un tenant, c'est-à-dire sur ses propres noms. Ce résidu
est accepté : il ne franchit aucune frontière de tenant.

\section{Politique de sélection des zones par adresse source}
\label{sec:vue}

Les deux options ouvertes par le document 03 sont de séparer physiquement les
points de résolution par tenant, ou d'appliquer une politique démontrée de
sélection des zones. La première option supposerait cinquante paires de
résolveurs, ce que le matériel disponible exclut. La seconde est donc retenue.

Son principe est direct. La relation \eqref{eq:prefixe} associe le tenant $n$ au
préfixe \ip{10.30.(63+n).0/24}. L'adresse source d'une requête détermine donc
sans ambiguïté le tenant émetteur, et par conséquent les deux seules zones
internes qu'il est autorisé à résoudre. En notant $\mathcal{A}$ la réunion des
préfixes d'administration et $o_3(a)$ le troisième octet de l'adresse $a$,
l'étiquette de vue s'écrit
\begin{equation}
  \tau(a) =
  \begin{cases}
    1 & \text{si } a \in \mathcal{A},\\[3pt]
    1000 + \bigl(o_3(a) - 63\bigr)
      & \text{si } a \in \texttt{10.30.64.0/18}
        \text{ et } 64 \leqslant o_3(a) \leqslant 113,\\[3pt]
    999 & \text{sinon,}
  \end{cases}
  \label{eq:tau}
\end{equation}
et l'ensemble des zones internes résolubles par un client d'étiquette
$\tau$ vaut
\begin{equation}
  \mathcal{Z}(\tau) =
  \begin{cases}
    \text{toutes les zones internes} & \text{si } \tau = 1,\\[3pt]
    \bigl\{\, \mathcal{D}(\tau-1000),\ \mathcal{R}(\tau-1000) \,\bigr\}
      & \text{si } \tau > 1000,\\[3pt]
    \varnothing & \text{si } \tau = 999 .
  \end{cases}
  \label{eq:zones}
\end{equation}

Toute demande portant sur l'espace interne et sortant de $\mathcal{Z}(\tau)$
reçoit un code \opt{REFUSED}, qui ne révèle pas l'existence du nom demandé,
contrairement à une réponse \opt{NXDOMAIN}.

\begin{encadredecision}
\titreencadredec{Point de conception critique : où placer le contrôle}
Une vérification placée dans le seul point d'entrée de résolution serait
insuffisante. Le résolveur consulte son cache paquet avant d'engager la
résolution : une réponse construite pour un tenant pourrait être restituée telle
quelle à un autre, et la politique serait contournée par le cache lui-même.

\smallskip
La fonction \opt{gettag()} est évaluée avant toute consultation de cache et sa
valeur de retour entre dans la clé du cache paquet. Le contrôle est donc
construit en deux temps : \opt{gettag()} calcule l'étiquette $\tau(a)$ de
l'équation \eqref{eq:tau}, puis \opt{preresolve()} applique la règle
\eqref{eq:zones} en s'appuyant sur cette étiquette. Deux vues distinctes ne
peuvent alors jamais partager une entrée de cache.
\end{encadredecision}

\begin{figure}[H]
\centering
\figChemin
\caption{Chemin d'une requête issue d'une machine tenant et points
  d'application de la politique}
\label{fig:chemin}
\end{figure}

\begin{table}[H]
\centering
\small
\caption{Vues définies par la politique de sélection des zones}\label{tab:vues}
\begin{tabular}{@{}C{0.12}R{0.41}R{0.47}@{}}
\hautfilet
\textbf{$\tau$} & \textbf{Adresses source} & \textbf{Portée de résolution}\\
\medfilet
1 & Plan de gestion \ip{10.30.0.0/24}, VNet Services \ip{10.30.4.0/24},
  Monitoring \ip{10.30.5.0/24}, pool VPN administrateur \ip{10.30.16.32/28} &
  Vue administration : résolution de l'ensemble de l'espace interne et de
  l'espace public.\\
$1000+n$ & Préfixe du tenant $n$, soit \ip{10.30.(63+n).0/24} &
  Zone directe $\mathcal{D}(n)$, zone inverse $\mathcal{R}(n)$, et espace
  public.\\
999 & Toute autre source acceptée par \opt{allow-from}, notamment le pool VPN
  utilisateur \ip{10.30.16.16/28} &
  Espace public uniquement ; tout nom interne reçoit \opt{REFUSED}.\\
aucune & Source hors \opt{allow-from} &
  La requête n'est pas traitée par le résolveur.\\
\basfilet
\end{tabular}
\end{table}

\section{Filtrage réseau et matrice de flux}
\label{sec:flux}

La politique applicative ne remplace pas le filtrage. Les règles suivantes
complètent les politiques de référence du document 02 et sont appliquées sur les
passerelles pfSense des VNets concernés, conformément aux interfaces IF-INT-05
et IF-INT-07.

\begin{longtable}{@{}R{0.21}R{0.17}R{0.13}R{0.16}R{0.33}@{}}
\caption{Matrice de flux du service DNS}\label{tab:matrice}\\
\hautfilet
\textbf{Source} & \textbf{Destination} & \textbf{Service} &
\textbf{Décision} & \textbf{Motif}\\
\medfilet
\endfirsthead
\hautfilet
\textbf{Source} & \textbf{Destination} & \textbf{Service} &
\textbf{Décision} & \textbf{Motif}\\
\medfilet
\endhead
\basfilet
\endfoot
VNet tenant \ip{10.30.(63+n).0/24} & \ip{10.30.4.12}, \ip{10.30.4.13} &
  UDP et TCP 53 & Autoriser & Seul point de résolution offert aux tenants\\
VNet tenant & \ip{10.30.4.10}, \ip{10.30.4.11} & tout &
  Refuser et journaliser & Les autoritatifs ne sont pas exposés aux tenants\\
VNet tenant & \ip{10.30.0.0/24} & tout & Refuser &
  Règle existante du document 02, rappelée ici\\
\ip{10.30.4.12}, \ip{10.30.4.13} & \ip{10.30.4.10}, \ip{10.30.4.11} &
  UDP et TCP 53 & Autoriser & Redirection de zone interne\\
\ip{10.30.4.12}, \ip{10.30.4.13} & Internet & UDP et TCP 53 &
  Autoriser via PF-WAN & Récursion sur l'espace public, journalisée\\
\ip{10.30.4.11} (ns2) & \ip{10.30.4.10} (ns1) & TCP 53 & Autoriser &
  Transfert de zone signé TSIG\\
\ip{10.30.4.10} (ns1) & \ip{10.30.4.11} (ns2) & UDP 53 & Autoriser &
  Message NOTIFY\\
\ip{10.30.4.41} (agent) & \ip{10.30.4.10} & TCP 8443 & Autoriser &
  API autoritative derrière terminaison TLS\\
\ip{10.30.4.41} (agent) & ns1, ns2, rec1, rec2 & UDP et TCP 53 & Autoriser &
  Contrôle de publication sur les quatre instances\\
\ip{10.30.0.0/24}, \ip{10.30.16.32/28} & \ip{10.30.4.12}, \ip{10.30.4.13} &
  UDP et TCP 53 & Autoriser & Vue administration\\
\ip{10.30.16.16/28} (VPN utilisateur) & \ip{10.30.4.12}, \ip{10.30.4.13} &
  UDP et TCP 53 & Autoriser &
  Admis par \opt{allow-from}, mais limité à l'espace public par la vue 999
  (section~\ref{sec:vpn})\\
\ip{10.30.5.10} (Monitoring) & ns1, ns2, rec1, rec2 & TCP 8443 & Autoriser &
  Lecture des métriques, initiée par Monitoring\\
ns1, ns2, rec1, rec2 & \ip{10.30.5.0/24} & tout & Refuser &
  Aucune initiation vers Monitoring, conformément à IF-NS-08\\
Toute autre source & ns1, ns2, rec1, rec2 & tout & Refuser &
  Refus par défaut\\
\end{longtable}

Deux points méritent attention. Le refus du trafic des tenants vers les
autoritatifs doit être journalisé : une hausse de ce compteur signale soit une
erreur de distribution des résolveurs, soit une tentative délibérée. Par
ailleurs, le port 53 doit être ouvert en TCP autant qu'en UDP, faute de quoi
toute réponse tronquée deviendrait irrécupérable, avec un symptôme limité à
certaines requêtes seulement.

\section{Cas particulier des profils VPN}
\label{sec:vpn}

L'exigence \exig{net-fr-017} demande que le DNS soit accessible selon le profil
VPN. Le plan d'adressage distingue trois pools : administrateur en
\ip{10.30.16.32/28}, utilisateur en \ip{10.30.16.16/28} et site à site en
\ip{10.30.16.48/28}, les deux derniers étant classés en extension.

Le pool administrateur ne pose pas de difficulté : il est traité comme le plan
de gestion et reçoit la vue administration, sous réserve de l'autorisation
délivrée par IAM au moment de l'ouverture de session.

Le pool utilisateur pose en revanche un problème de conception qu'il faut
énoncer plutôt que masquer. Une politique fondée sur l'adresse source ne peut
pas distinguer deux utilisateurs de tenants différents qui reçoivent leur
adresse d'un même pool partagé de quatorze adresses. Trois réponses sont
possibles.

\begin{enumerate}
  \item Attribuer une sous-plage distincte par tenant, ce que la taille du pool
    actuel ne permet pas pour cinquante tenants.
  \item Faire porter l'identité du tenant par le concentrateur VPN, qui
    attribuerait une adresse stable par session et transmettrait la
    correspondance au résolveur, au titre de l'interface IF-EXT-07.
  \item Refuser toute résolution interne au pool utilisateur et médier les accès
    par le Backend, conformément à IF-EXT-06.
\end{enumerate}

La troisième réponse est retenue pour le périmètre minimal, parce qu'elle est la
seule applicable sans donnée supplémentaire et parce qu'elle est cohérente avec
le statut d'extension du pool utilisateur. L'étiquette $\tau = 999$ de
l'équation \eqref{eq:tau} matérialise ce refus. La deuxième réponse constitue
l'évolution naturelle dès que la correspondance session vers tenant sera
fournie.

% ================== Chapitre 4 : configurations de reference ================
\chapter{Configurations de référence}
\label{ch:configurations}

Les configurations de ce chapitre sont complètes et cohérentes entre elles.
Elles correspondent à la branche 4.9 de PowerDNS Authoritative et à la branche 5
de PowerDNS Recursor. Deux précautions de version s'imposent avant application :
les directives \opt{primary} et \opt{secondary} remplacent les anciennes
directives \opt{master} et \opt{slave} depuis la version 4.5, et la branche 5 du
résolveur introduit un format YAML équivalent au format classique employé ici.
Le format attendu par la version effectivement installée doit être vérifié.

\section{Serveurs autoritatifs ns1 et ns2}
\label{sec:autoritatifs}

\begin{lstlisting}[language=pdnsconf,caption={Configuration du serveur
  autoritatif primaire ns1},label=lst:ns1]
# /etc/powerdns/pdns.conf
# GANDAL T005 : ns1, serveur autoritatif primaire, 10.30.4.10, hote H1

# --- stockage : volume faible, replication assuree par transfert de zone
launch=gsqlite3
gsqlite3-database=/var/lib/powerdns/pdns.sqlite3

# --- ecoute restreinte a l'adresse de service du VNet Services
local-address=10.30.4.10
local-port=53

# --- role
primary=yes
secondary=no

# --- replication : ns2 est le seul destinataire et le seul demandeur autorise
allow-axfr-ips=10.30.4.11/32
only-notify=10.30.4.11
also-notify=10.30.4.11

# --- valeurs par defaut des zones creees (cf. politique de duree de vie)
default-ttl=3600
default-soa-content=ns1.gandal.internal hostmaster.gandal.internal 0 3600 900 604800 300

# --- transport : rester sous la MTU utile de 1450 octets imposee par VXLAN
udp-truncation-threshold=1232
max-tcp-connections=50

# --- discretion et isolement
version-string=anonymous
security-poll-suffix=

# --- API locale ; la terminaison TLS est assuree par le mandataire (4.5)
webserver=yes
webserver-address=127.0.0.1
webserver-port=8081
webserver-allow-from=127.0.0.1/32
api=yes
api-key=$PDNS_API_KEY_NS1

# --- journalisation
loglevel=4
log-dns-queries=no
log-dns-details=yes

# --- execution sous systemd
setuid=pdns
setgid=pdns
daemon=no
guardian=no
\end{lstlisting}

Quatre directives appellent un commentaire. La journalisation des requêtes est
désactivée sur l'autoritatif : ses clients légitimes sont au nombre de quatre et
l'identité réelle du demandeur n'y est pas visible, puisque les requêtes
arrivent par le résolveur. La journalisation utile est placée sur le résolveur,
où l'adresse source est celle du client. La directive \opt{security-poll-suffix}
vidée désactive la vérification périodique de version auprès d'un domaine
externe, requête inutile et indésirable dans un environnement fermé. La
directive \opt{version-string} masque la version du logiciel dans la réponse à
\zone{version.bind}. Enfin, la clé d'API n'est pas écrite dans le fichier : elle
est injectée dans l'environnement du service au démarrage, depuis le magasin de
secrets. À partir de la version 4.7, la directive \opt{api-key} accepte également
une empreinte produite par la commande \opt{pdnsutil hash-password}, ce qui évite
toute forme en clair.

\begin{lstlisting}[language=pdnsconf,caption={Écarts de configuration du serveur
  autoritatif secondaire ns2},label=lst:ns2]
# /etc/powerdns/pdns.conf
# GANDAL T005 : ns2, serveur autoritatif secondaire, 10.30.4.11, hote H3
# Seules les directives qui different de ns1 sont reproduites.

local-address=10.30.4.11

primary=no
secondary=yes

# --- ns1 est la seule source de NOTIFY et de transfert acceptee
allow-notify-from=10.30.4.10/32
xfr-cycle-interval=60
xfr-max-received-mbytes=50

# --- aucun transfert sortant : ns2 ne sert de source a personne
allow-axfr-ips=

api-key=$PDNS_API_KEY_NS2
\end{lstlisting}

La directive \opt{allow-axfr-ips} laissée vide ferme toute possibilité de
transfert depuis ns2. L'API de ns2 reste active pour la supervision et
l'inspection locale, mais le mandataire placé devant elle n'autorise que la
méthode \opt{GET}, de sorte que le propriétaire d'écriture reste unique au sens
de la section~\ref{sec:proprietaire}.

\section{Création des zones et des clés de transfert}
\label{sec:zones}

La clé de transfert est unique parce que l'unique transfert du dispositif va de
ns1 vers ns2. Le document 03 évoque des clés TSIG limitées par zone : cette
granularité devient nécessaire dans la variante où les mises à jour dynamiques
du RFC 2136 seraient employées, chaque émetteur ne devant alors pouvoir modifier
que sa propre zone. Cette variante n'est pas retenue, comme l'explique la
section~\ref{sec:proprietaire}.

\begin{lstlisting}[caption={Création de l'espace de nommage sur ns1},
  label=lst:zones]
# --- sur ns1 : generation de la cle de transfert
pdnsutil generate-tsig-key gandal-xfr hmac-sha256
# la valeur produite est deposee dans le magasin de secrets, jamais dans Git

# --- zone racine interne et identites des serveurs
pdnsutil create-zone gandal.internal ns1.gandal.internal
pdnsutil add-record  gandal.internal @    NS 3600 ns2.gandal.internal.
pdnsutil add-record  gandal.internal ns1  A  3600 10.30.4.10
pdnsutil add-record  gandal.internal ns2  A  3600 10.30.4.11
pdnsutil add-record  gandal.internal rec1 A  3600 10.30.4.12
pdnsutil add-record  gandal.internal rec2 A  3600 10.30.4.13

# --- zone d'infrastructure, puis sa delegation depuis la zone parente
pdnsutil create-zone infra.gandal.internal ns1.gandal.internal
pdnsutil add-record  infra.gandal.internal @ NS 3600 ns2.gandal.internal.
pdnsutil add-record  gandal.internal infra NS 3600 ns1.gandal.internal.
pdnsutil add-record  gandal.internal infra NS 3600 ns2.gandal.internal.

# --- inventaire d'infrastructure, extrait
pdnsutil add-record infra.gandal.internal h1      A 3600 10.30.0.11
pdnsutil add-record infra.gandal.internal h2      A 3600 10.30.0.12
pdnsutil add-record infra.gandal.internal h3      A 3600 10.30.0.13
pdnsutil add-record infra.gandal.internal xo      A 3600 10.30.0.20
pdnsutil add-record infra.gandal.internal bastion A 3600 10.30.0.21
pdnsutil add-record infra.gandal.internal kea1    A 3600 10.30.4.20
pdnsutil add-record infra.gandal.internal kea2    A 3600 10.30.4.21
pdnsutil add-record infra.gandal.internal ntp1    A 3600 10.30.4.30
pdnsutil add-record infra.gandal.internal ntp2    A 3600 10.30.4.31
pdnsutil add-record infra.gandal.internal ipam    A 3600 10.30.4.40
pdnsutil add-record infra.gandal.internal agent   A 3600 10.30.4.41

# --- zones inverses d'infrastructure
for O in 0 1 2 4 5 6 7 8 10 16 17 ; do
  pdnsutil create-zone "${O}.30.10.in-addr.arpa" ns1.gandal.internal
  pdnsutil add-record  "${O}.30.10.in-addr.arpa" @ NS 3600 ns2.gandal.internal.
done

# --- cinquante tenants : zone directe, zone inverse et delegation
for N in $(seq 1 50) ; do
  T=$(printf 'tenant-%03d' "$N")
  O=$(( 63 + N ))
  pdnsutil create-zone "${T}.gandal.internal" ns1.gandal.internal
  pdnsutil add-record  "${T}.gandal.internal" @  NS 3600 ns2.gandal.internal.
  pdnsutil add-record  "${T}.gandal.internal" gw A  3600 "10.30.${O}.1"
  pdnsutil create-zone "${O}.30.10.in-addr.arpa" ns1.gandal.internal
  pdnsutil add-record  "${O}.30.10.in-addr.arpa" @ NS  3600 ns2.gandal.internal.
  pdnsutil add-record  "${O}.30.10.in-addr.arpa" 1 PTR 3600 "gw.${T}.gandal.internal."
  pdnsutil add-record  gandal.internal "$T" NS 3600 ns1.gandal.internal.
  pdnsutil add-record  gandal.internal "$T" NS 3600 ns2.gandal.internal.
done

# --- transfert signe et increment automatique du numero de serie par l'API
for Z in $(pdnsutil list-all-zones) ; do
  pdnsutil set-meta "$Z" TSIG-ALLOW-AXFR gandal-xfr
  pdnsutil set-meta "$Z" SOA-EDIT-API    DEFAULT
done

# --- controle de coherence avant mise en service
pdnsutil rectify-all-zones
pdnsutil check-all-zones
\end{lstlisting}

\begin{lstlisting}[caption={Déclaration des zones secondaires sur ns2},
  label=lst:ns2zones]
# --- sur ns2 : import de la meme cle, puis declaration des zones secondaires
pdnsutil import-tsig-key gandal-xfr hmac-sha256 "$TSIG_GANDAL_XFR"

for Z in $(cat /etc/gandal/zones.list) ; do
  pdnsutil create-secondary-zone "$Z" 10.30.4.10
  pdnsutil activate-tsig-key     "$Z" gandal-xfr secondary
done
\end{lstlisting}

La métadonnée \opt{SOA-EDIT-API} fixée à \opt{DEFAULT} fait incrémenter le
numéro de série à chaque modification reçue par l'API. C'est cette
incrémentation qui déclenche le message \opt{NOTIFY} et le transfert vers ns2 ;
sans elle, les deux serveurs divergeraient silencieusement.

\section{Résolveurs récursifs rec1 et rec2}
\label{sec:recursif}

\begin{lstlisting}[language=pdnsconf,caption={Configuration du résolveur
  récursif rec1},label=lst:rec1]
# /etc/powerdns/recursor.conf
# GANDAL T005 : rec1, resolveur recursif, 10.30.4.12, hote H2
# rec2 est identique, avec local-address=10.30.4.13 sur l'hote H1

local-address=10.30.4.12
local-port=53
threads=2
pdns-distributes-queries=no

# --- clients autorises : strictement les prefixes du plan 02
allow-from=10.30.0.0/24, 10.30.4.0/24, 10.30.5.0/24, 10.30.6.0/24,
           10.30.7.0/24, 10.30.8.0/24, 10.30.16.16/28, 10.30.16.32/28,
           10.30.64.0/18

# --- espace interne : redirection vers les serveurs autoritatifs
forward-zones=gandal.internal=10.30.4.10;10.30.4.11
forward-zones+=30.10.in-addr.arpa=10.30.4.10;10.30.4.11

# Sans cette directive, le resolveur servirait localement les zones inverses
# privees du RFC 6303 et court-circuiterait la redirection ci-dessus.
serve-rfc1918=no

# --- validation DNSSEC de l'espace public, exception traitee dans le Lua
dnssec=validate
lua-config-file=/etc/powerdns/recursor-config.lua

# --- politique de selection des zones (section 4.4)
lua-dns-script=/etc/powerdns/views.lua

# --- transport : rester sous la MTU utile de 1450 octets imposee par VXLAN
edns-outgoing-bufsize=1232
udp-truncation-threshold=1232
qname-minimization=yes

# --- cache dimensionne pour 4 GiB de memoire
max-cache-entries=500000
max-packetcache-entries=250000
max-cache-ttl=86400
max-negative-ttl=300

# --- discretion et isolement
version-string=anonymous
security-poll-suffix=
export-etc-hosts=no

# --- point de metriques, expose par le mandataire (section 4.5)
webserver=yes
webserver-address=127.0.0.1
webserver-port=8082
webserver-allow-from=127.0.0.1/32
api-key=$PDNS_REC_API_KEY

loglevel=4
\end{lstlisting}

La directive \opt{serve-rfc1918} mérite d'être soulignée. Par défaut, le
résolveur se déclare autoritatif et vide pour les zones inverses des espaces
privés, afin d'éviter des requêtes inutiles vers la racine. Dans GANDAL, ce
comportement empêcherait toute résolution inverse interne. Il est donc désactivé
et remplacé par une redirection explicite de \zone{30.10.in-addr.arpa} vers les
autoritatifs.

La directive \opt{export-etc-hosts} est maintenue à la valeur \opt{no} : le
fichier \opt{hosts} local des résolveurs contient le socle de secours décrit à la
section~\ref{sec:amorcage}, qui n'a pas vocation à être publié et qui pourrait
entrer en contradiction avec les zones autoritatives.

\begin{lstlisting}[language=luapdns,caption={Configuration Lua du résolveur :
  exception de validation et journalisation},label=lst:reclua]
-- /etc/powerdns/recursor-config.lua
-- Charge par la directive lua-config-file.

-- Le domaine de premier niveau "internal" est reserve a l'usage prive et n'est
-- pas delegue dans la racine publique. Avec dnssec=validate, le resolveur
-- obtiendrait de la racine une preuve signee de son inexistence et declarerait
-- toute reponse interne invalide. L'exception locale ci-dessous suspend la
-- validation sous cette branche, et sous elle seulement.
addNTA("internal.", "espace de nommage interne GANDAL, non delegue et non signe")

-- Aucune exception n'est necessaire pour 10.in-addr.arpa : cette branche est
-- deleguee sans enregistrement DS dans l'arbre public, donc deja non signee.

-- Journalisation structuree vers un collecteur local du VNet Services.
-- Monitoring lit ensuite les fichiers produits en pull (IF-NS-08).
dnstapFrameStreamServer({"/var/run/pdns-recursor/dnstap.sock"})
\end{lstlisting}

\section{Script de politique de vue}
\label{sec:script}

Le script ci-dessous met en \oe uvre les équations \eqref{eq:tau} et
\eqref{eq:zones}. Il est volontairement court et sans accès externe : il
s'exécute pour chaque requête et toute opération coûteuse s'y traduirait
directement en latence de résolution.

\begin{lstlisting}[language=luapdns,caption={Politique de sélection des zones
  appliquée sur rec1 et rec2},label=lst:views]
-- /etc/powerdns/views.lua
-- Politique de selection des zones par adresse source.
-- Exigences couvertes : NET-FR-008, IF-NS-03 ; scenario de test S04.

local ZONE_RACINE = "gandal.internal"
local SUFFIXE_INV = "30.10.in-addr.arpa"

local TAG_ADMIN = 1     -- gestion, services, monitoring, VPN administrateur
local TAG_REFUS = 999   -- source autorisee mais sans vue interne

local reseauxAdmin = newNMG()
reseauxAdmin:addMask("10.30.0.0/24")
reseauxAdmin:addMask("10.30.4.0/24")
reseauxAdmin:addMask("10.30.5.0/24")
reseauxAdmin:addMask("10.30.16.32/28")

local reseauxTenant = newNMG()
reseauxTenant:addMask("10.30.64.0/18")

-- Plan 02 : le tenant n occupe 10.30.(63+n).0/24, pour 1 <= n <= 50.
local function numeroTenant(adresse)
  local o3 = tonumber(adresse:toString():match("^10%.30%.(%d+)%.%d+$"))
  if o3 == nil or o3 < 64 or o3 > 113 then return nil end
  return o3 - 63
end

-- gettag() est evalue avant toute consultation de cache et sa valeur entre
-- dans la cle du cache paquet : deux vues ne partagent jamais une reponse.
function gettag(remote, ednssubnet, localip, qname, qtype, ednsoptions, tcp)
  if reseauxAdmin:match(remote) then return TAG_ADMIN end
  if reseauxTenant:match(remote) then
    local n = numeroTenant(remote)
    if n ~= nil then return 1000 + n end
  end
  return TAG_REFUS
end

local function sousDomaine(nom, zone)
  return nom == zone or nom:sub(-(#zone + 1)) == ("." .. zone)
end

function preresolve(dq)
  local nom     = dq.qname:toStringNoDot():lower()
  local interne = sousDomaine(nom, ZONE_RACINE)
  local inverse = sousDomaine(nom, SUFFIXE_INV)

  -- Les noms publics suivent la recursion ordinaire, quelle que soit la vue.
  if not interne and not inverse then return false end

  if dq.tag == TAG_ADMIN then return false end

  if dq.tag > 1000 then
    local n           = dq.tag - 1000
    local zoneDirecte = string.format("tenant-%03d.%s", n, ZONE_RACINE)
    local zoneInverse = string.format("%d.%s", n + 63, SUFFIXE_INV)
    if sousDomaine(nom, zoneDirecte) or sousDomaine(nom, zoneInverse) then
      return false
    end
  end

  -- REFUSED plutot que NXDOMAIN : le refus ne revele pas l'existence du nom.
  dq.rcode = pdns.REFUSED
  return true
end
\end{lstlisting}

Le script refuse également l'apex \zone{gandal.internal} aux sources tenant, ce
qui empêche l'énumération des délégations et donc la découverte de la liste des
tenants. Il laisse en revanche passer sans traitement toute requête portant sur
un nom public, afin de ne pas pénaliser le cas le plus fréquent.

\section{Exposition de l'API et des métriques}
\label{sec:api}

Le serveur HTTP intégré à PowerDNS ne fournit pas de terminaison TLS.
L'exigence \exig{net-fr-016}, qui demande des API administratives chiffrées, est
donc satisfaite en liant ce serveur à la boucle locale et en plaçant devant lui
un mandataire qui assure le chiffrement et le contrôle d'accès par adresse.

\begin{lstlisting}[caption={Terminaison TLS et contrôle d'accès devant l'API
  autoritative},label=lst:nginx]
# /etc/nginx/conf.d/pdns-api.conf  sur ns1 (10.30.4.10)

server {
    listen 10.30.4.10:8443 ssl;
    server_name ns1.gandal.internal;

    ssl_certificate     /etc/ssl/gandal/ns1.crt;   # PKI interne, IF-EXT-01
    ssl_certificate_key /etc/ssl/gandal/ns1.key;
    ssl_protocols       TLSv1.2 TLSv1.3;

    # Seul l'agent reseau publie des enregistrements.
    location /api/ {
        allow 10.30.4.41/32;
        deny  all;
        proxy_pass http://127.0.0.1:8081;
    }

    # Seul Monitoring lit les metriques, et il en est l'initiateur.
    location = /metrics {
        allow 10.30.5.10/32;
        deny  all;
        proxy_pass http://127.0.0.1:8081/metrics;
    }

    location / { return 404; }
}

# Sur ns2, le bloc /api/ est complete par la restriction suivante, qui
# garantit l'unicite du proprietaire d'ecriture :
#     limit_except GET { deny all; }
#
# Sur rec1 et rec2, seul le bloc /metrics est publie, vers 127.0.0.1:8082.
\end{lstlisting}

Ce montage répond également à l'interface IF-NS-08 : les quatre instances
n'ouvrent aucune connexion vers Monitoring, qui vient lire leur point de
métriques. La limitation de débit en réponse n'est pas activée sur les serveurs
autoritatifs, et cette absence est délibérée : PowerDNS Authoritative n'intègre
pas cette fonction, et le risque d'amplification qu'elle traite suppose une
exposition publique que la matrice de flux interdit. Si cette exposition devait
changer, un répartiteur DNS spécialisé devrait être placé en amont.

% ================== Chapitre 5 : cycle de vie et exploitation ===============
\chapter{Cycle de vie des enregistrements et exploitation}
\label{ch:cycle}

\section{Propriétaire d'écriture unique et séquence de publication}
\label{sec:proprietaire}

Le document 03 pose une règle explicite : ne choisir qu'un propriétaire
d'écriture par enregistrement. Il ouvre deux pistes, la mise à jour dynamique
issue de Kea au moyen du RFC 2136, ou la mise à jour par l'agent au travers de
l'API autoritative.

\begin{table}[H]
\centering
\small
\caption{Arbitrage du mécanisme de mise à jour}\label{tab:arbitrage}
\begin{tabular}{@{}R{0.22}R{0.30}R{0.33}R{0.15}@{}}
\hautfilet
\textbf{Option} & \textbf{Apport} & \textbf{Limite} & \textbf{Décision}\\
\medfilet
Mise à jour par l'agent via l'API autoritative &
  Un seul émetteur pour toutes les zones. La modification porte sur un RRset
  complet, ce qui rend l'opération naturellement convergente. Le retour d'API
  fournit un état exploitable par la machine d'états et corrélable à un
  \opt{operation\_id}. &
  Ajoute une dépendance à la disponibilité de l'agent pour toute publication. &
  \textbf{Retenue}\\
Mise à jour dynamique RFC 2136 déclenchée par Kea &
  Rapproche la publication de l'attribution du bail et supprime un
  intermédiaire. &
  Introduit un second émetteur, donc un conflit possible de propriété sur les
  mêmes enregistrements. Exige une clé TSIG par zone, soit cent treize clés à
  gérer, et une compatibilité du moteur de stockage qui reste à vérifier. &
  Alternative documentée\\
\basfilet
\end{tabular}
\end{table}

\begin{figure}[H]
\centering
\figSequence
\caption{Séquence de publication d'un enregistrement et branche dégradée}
\label{fig:sequence}
\end{figure}

\label{sec:publication}
Les deux appels ci-dessous constituent le contrat que l'agent doit respecter.
Ils sont donnés pour une machine du tenant 7 recevant l'adresse
\ip{10.30.70.37}, qui appartient bien au pool dynamique \ip{.10} à \ip{.199}
défini par le document 02.

\begin{lstlisting}[caption={Publication conjointe de l'enregistrement direct et
  de l'enregistrement inverse},label=lst:publication]
PATCH /api/v1/servers/localhost/zones/tenant-007.gandal.internal.  HTTP/1.1
Host: ns1.gandal.internal:8443
X-API-Key: <cle injectee depuis le magasin de secrets>
Content-Type: application/json

{
  "rrsets": [{
    "name": "vm-0042.tenant-007.gandal.internal.",
    "type": "A",
    "ttl": 300,
    "changetype": "REPLACE",
    "records": [ { "content": "10.30.70.37", "disabled": false } ],
    "comments": [ { "account": "agent",
                    "content": "operation_id=7f3c..., lease_id=..." } ]
  }]
}

PATCH /api/v1/servers/localhost/zones/70.30.10.in-addr.arpa.  HTTP/1.1

{
  "rrsets": [{
    "name": "37.70.30.10.in-addr.arpa.",
    "type": "PTR",
    "ttl": 300,
    "changetype": "REPLACE",
    "records": [ { "content": "vm-0042.tenant-007.gandal.internal.",
                   "disabled": false } ]
  }]
}
\end{lstlisting}

La publication n'est pas considérée comme acquise à la réception du code 204.
L'agent interroge ensuite ns1, ns2, rec1 et rec2 pour le nom publié et ne fait
progresser l'opération vers l'état \opt{SERVICES\_READY} qu'après concordance
des quatre réponses. Ce contrôle détecte en particulier un transfert de zone
bloqué, défaut qui resterait sinon invisible jusqu'à la première panne de ns1.

\section{Idempotence, rejeu et réconciliation}
\label{sec:idempotence}

L'exigence \exig{net-nfr-006} demande un rejeu déterministe et idempotent. Le
choix d'un changement de type \opt{REPLACE} portant sur un RRset complet y
répond par construction : en notant $\Phi$ l'application d'une publication sur
l'état des zones, $\Phi \circ \Phi = \Phi$, c'est-à-dire que l'état final ne
dépend pas du nombre d'exécutions. Une précaution reste nécessaire.

\begin{encadre}
\titreencadre{Précaution sur le rejeu}
Une modification par l'API incrémente le numéro de série de la zone, y compris
lorsque le contenu soumis est identique à celui déjà publié. Un rejeu
systématique provoquerait donc autant de transferts de zone inutiles vers ns2.

\smallskip
L'agent doit par conséquent lire le RRset avant de l'écrire et n'émettre la
modification que si le contenu diffère. Le rejeu reste correct dans tous les
cas, mais il devient également économe.
\end{encadre}

Une réconciliation périodique compare les allocations enregistrées dans l'IPAM
aux enregistrements effectivement publiés, dans les deux sens. Elle signale les
divergences sans les corriger automatiquement : une correction automatique sur
un état mal compris est le moyen le plus rapide de propager une erreur à
l'ensemble des zones. Deux familles de divergence sont attendues,
l'enregistrement publié sans allocation correspondante, trace d'une suppression
incomplète, et l'allocation active sans enregistrement, trace d'une publication
interrompue.

\section{Suppression, quarantaine et mode dégradé}
\label{sec:suppression}

L'ordre de suppression est l'inverse de l'ordre de publication et il n'est pas
interchangeable. Le document 03 impose d'attendre le retrait des interfaces
virtuelles et la fin des références avant libération, et précise qu'une panne
DNS ne libère pas une adresse encore utilisée.

\begin{enumerate}
  \item Retrait de l'interface virtuelle de la machine et fin du bail.
  \item Suppression des enregistrements A et PTR par un changement de type
    \opt{DELETE} sur les deux zones concernées.
  \item Attente de la durée de quarantaine $Q = 900$ secondes, qui satisfait
    l'inégalité \eqref{eq:ttl}.
  \item Retour de l'adresse au pool par l'IPAM.
\end{enumerate}

Si la suppression des enregistrements échoue, l'adresse n'est pas rendue au pool
et l'opération reste dans l'état \opt{DELETING}, observable. Le comportement
symétrique s'applique à la publication : une panne de l'API ou des serveurs
autoritatifs laisse l'opération en cours, signalée par une alerte, sans que
l'adresse déjà attribuée soit considérée comme libre. Le trafic des machines
déjà configurées n'est pas interrompu : seules les opérations nouvelles sont
suspendues. Cette distinction est précisément ce que le test \opt{G04} doit
établir.

\section{Amorçage, ordre de reprise et adresses de secours}
\label{sec:amorcage}

Le document Underlay insiste sur la rupture des cycles de démarrage : si le
serveur DNS et le bastion dépendent uniquement d'un ensemble dont l'accès
suppose ces mêmes services, une panne générale rend la restauration difficile.
Le socle minimal ci-dessous doit donc être présent sur les trois hôtes, sur XO,
sur le bastion et sur les machines de services.

\begin{lstlisting}[language=pdnsconf,caption={Adresses de secours à conserver
  hors du service DNS},label=lst:hosts]
# /etc/hosts  : socle de reprise GANDAL
# Permet d'administrer et de reconstruire l'infrastructure sans dependre
# du service DNS lui-meme. Non publie : voir export-etc-hosts=no.

127.0.0.1    localhost
10.30.0.11   h1.infra.gandal.internal         h1
10.30.0.12   h2.infra.gandal.internal         h2
10.30.0.13   h3.infra.gandal.internal         h3
10.30.0.20   xo.infra.gandal.internal         xo
10.30.0.21   bastion.infra.gandal.internal    bastion
10.30.4.10   ns1.gandal.internal              ns1
10.30.4.11   ns2.gandal.internal              ns2
10.30.4.12   rec1.gandal.internal             rec1
10.30.4.13   rec2.gandal.internal             rec2
10.30.4.20   kea1.infra.gandal.internal       kea1
10.30.4.21   kea2.infra.gandal.internal       kea2
10.30.4.30   ntp1.infra.gandal.internal       ntp1
10.30.4.31   ntp2.infra.gandal.internal       ntp2
10.30.4.40   ipam.infra.gandal.internal       ipam
10.30.4.41   agent.infra.gandal.internal      agent
\end{lstlisting}

\begin{longtable}{@{}C{0.09}R{0.24}R{0.15}R{0.52}@{}}
\caption{Ordre de reprise après indisponibilité étendue}\label{tab:reprise}\\
\hautfilet
\textbf{Rang} & \textbf{Élément remis en service} & \textbf{Dépend de} &
\textbf{Remarque}\\
\medfilet
\endfirsthead
\hautfilet
\textbf{Rang} & \textbf{Élément remis en service} & \textbf{Dépend de} &
\textbf{Remarque}\\
\medfilet
\endhead
\basfilet
\endfoot
1 & Transport et hôtes & Aucune &
  Les hyperviseurs se joignent par adresse de gestion, sans résolution.\\
2 & Source de temps & Transport &
  La validation DNSSEC et la vérification des certificats exigent une horloge
  plausible. Au démarrage à froid, le client de temps doit être autorisé à
  corriger un écart important avant toute vérification cryptographique.\\
3 & Autoritatifs ns1 puis ns2 & Temps &
  Le fichier \opt{hosts} suffit à les atteindre.\\
4 & Résolveurs rec1 et rec2 & Autoritatifs &
  Leur propre résolution locale pointe sur la boucle locale, jamais sur leur
  homologue, afin d'éviter une dépendance croisée.\\
5 & IPAM puis agent & DNS &
  La réconciliation ne démarre qu'une fois le DNS joignable.\\
6 & DHCP & DNS et IPAM &
  Les options distribuées référencent les résolveurs.\\
7 & Orchestration et supervision & Ensemble des précédents &
  Dernière étape ; la supervision constate l'état sans conditionner la remise
  en service.\\
\end{longtable}

Le point sensible est le rang 2. La validation cryptographique du temps par NTS
suppose un certificat valide, dont la vérification suppose une horloge correcte,
qui suppose elle-même le service de temps. Le document 03 signale ce cas sans le
trancher. Du point de vue du DNS, la conséquence est simple : si l'horloge est
fausse, la validation DNSSEC rejette des réponses publiques pourtant légitimes,
ce qui se manifeste par des codes \opt{SERVFAIL} sur l'espace public alors que
l'espace interne répond normalement. Ce symptôme doit figurer dans la procédure
de diagnostic, car il oriente immédiatement vers l'horloge et non vers le DNS.

\section{Journalisation, métriques et alertes}
\label{sec:journaux}

La journalisation est placée sur les résolveurs, seuls points où l'adresse
source est celle du client réel. Elle est horodatée en temps universel
coordonné, conformément à \exig{net-fr-018}, et limitée aux éléments
nécessaires : horodatage, adresse source, nom demandé, type, code de réponse et
étiquette de vue. La durée de conservation et la responsabilité des données
relèvent de l'interface avec IAM et les responsables métier.

\begin{longtable}{@{}R{0.30}R{0.30}R{0.40}@{}}
\caption{Indicateurs et alertes du service DNS}\label{tab:alertes}\\
\hautfilet
\textbf{Indicateur} & \textbf{Seuil ou condition} & \textbf{Portée}\\
\medfilet
\endfirsthead
\hautfilet
\textbf{Indicateur} & \textbf{Seuil ou condition} & \textbf{Portée}\\
\medfilet
\endhead
\basfilet
\endfoot
Disponibilité de chaque instance & Absence de réponse à une sonde SOA pendant
  60 s & Majeure\\
Écart de numéro de série entre ns1 et ns2 & Écart persistant au-delà de deux
  cycles de transfert, soit 120 s &
  Majeure : la réplication est rompue et la redondance n'est plus effective\\
Taux de réponses \opt{REFUSED} par vue & Hausse significative par rapport à la
  référence &
  Mineure : politique trop stricte, mauvaise distribution des résolveurs, ou
  tentative d'accès\\
Taux de réponses \opt{SERVFAIL} & Hausse sur l'espace public &
  Majeure : dérive d'horloge, perte de la sortie externe ou échec de
  validation\\
Âge du dernier transfert réussi sur ns2 & Supérieur à deux fois
  \opt{refresh}, soit 7\,200 s & Majeure\\
Opérations de publication en échec & Toute opération bloquée au-delà de son
  délai & Majeure\\
Divergence IPAM et zones publiées & Toute divergence persistant après deux
  cycles de réconciliation & Mineure, à instruire\\
Tentatives de transfert de zone refusées & Toute occurrence &
  À instruire en sécurité\\
Absence de données de collecte & Interruption de la lecture en pull &
  Majeure : la panne de collecte doit alerter, et non effacer silencieusement
  les traces\\
\end{longtable}

\section{Sauvegarde, restauration et gestion des secrets}
\label{sec:sauvegarde}

L'objectif historique de reconstruction en moins de soixante minutes reste à
évaluer, et cette évaluation relève de T009. Ce dossier en fournit les éléments :
ce qui doit être sauvegardé, où, et dans quel ordre la restauration s'effectue.

\begin{lstlisting}[caption={Export et restauration de l'état autoritatif},
  label=lst:sauvegarde]
# --- export quotidien de l'etat autoritatif, depuis ns1
install -d -m 0750 /var/backups/gandal/dns
for Z in $(pdnsutil list-all-zones) ; do
  pdnsutil list-zone "$Z" > "/var/backups/gandal/dns/${Z}.zone"
done

# --- l'export est verse dans le depot versionne, sans aucun secret
#     Les cles TSIG et les cles d'API restent dans le magasin de secrets.

# --- restauration de ns1 sur une machine neuve
pdnsutil import-tsig-key gandal-xfr hmac-sha256 "$TSIG_GANDAL_XFR"
for F in /var/backups/gandal/dns/*.zone ; do
  Z=$(basename "$F" .zone)
  pdnsutil create-zone "$Z" ns1.gandal.internal
  pdnsutil load-zone   "$Z" "$F"
  pdnsutil set-meta    "$Z" TSIG-ALLOW-AXFR gandal-xfr
  pdnsutil set-meta    "$Z" SOA-EDIT-API    DEFAULT
done
pdnsutil rectify-all-zones
pdnsutil check-all-zones
\end{lstlisting}

\begin{table}[H]
\centering
\small
\caption{Gestion des secrets et des sauvegardes du service}\label{tab:secrets}
\begin{tabular}{@{}R{0.22}R{0.28}R{0.50}@{}}
\hautfilet
\textbf{Élément} & \textbf{Emplacement} & \textbf{Règle}\\
\medfilet
Clé TSIG de transfert & Magasin de secrets &
  Rotation annuelle ou sur incident. La rotation se fait en déclarant la
  nouvelle clé sur les deux serveurs avant de retirer l'ancienne.\\
Clé d'API de ns1 & Magasin de secrets, injectée dans l'environnement du
  service &
  Rotation lors de tout changement du porteur de l'agent. L'empreinte peut
  figurer en configuration à partir de la version 4.7.\\
Certificat TLS du mandataire & PKI interne, interface IF-EXT-01 &
  Renouvellement avant échéance ; la supervision du délai restant fait partie
  des indicateurs attendus.\\
Exports de zones & Dépôt versionné et copie hors du cluster &
  Une copie qui ne serait disponible que dans le cluster administré n'offre
  aucune reprise autonome après sa perte.\\
\basfilet
\end{tabular}
\end{table}

% ================== Chapitre 6 : controles de recette =======================
\chapter{Contrôles de recette de la tâche}
\label{ch:recette}

\section{Critères d'acceptation et commandes associées}

Les contrôles ci-dessous définissent ce que signifie exactement, pour T005, que
le service DNS est en place. Ils portent sur le service lui-même et non sur le
comportement d'ensemble du datacenter. Conformément à la règle du corpus, leur
statut initial est \statut{non exécuté} et aucun résultat n'est présenté comme
obtenu.

\begin{longtable}{@{}C{0.09}R{0.19}R{0.37}R{0.35}@{}}
\caption{Contrôles de recette de T005 ; statut initial de l'ensemble :
  \statut{non exécuté}}\label{tab:controles}\\
\hautfilet
\textbf{Code} & \textbf{Objet} & \textbf{Précondition et action} &
\textbf{Résultat attendu}\\
\medfilet
\endfirsthead
\hautfilet
\textbf{Code} & \textbf{Objet} & \textbf{Précondition et action} &
\textbf{Résultat attendu}\\
\medfilet
\endhead
\basfilet
\endfoot
C01 & Réponse autoritative directe &
  \texttt{dig +norecurse @10.30.4.10 ns1.gandal.internal A} &
  \opt{NOERROR}, indicateur \opt{aa} présent, réponse \ip{10.30.4.10}\\
C02 & Cohérence entre les deux autoritatifs &
  \texttt{dig +short @10.30.4.10 tenant-007.gandal.internal SOA} puis la même
  requête sur \ip{10.30.4.11} &
  Numéros de série identiques après réception du \opt{NOTIFY}\\
C03 & Résolution inverse d'une machine tenant &
  \texttt{dig +short @10.30.4.12 -x 10.30.70.37} &
  \zone{vm-0042.tenant-007.\allowbreak gandal.internal.}\\
C04 & Transfert de zone refusé à un tiers &
  \texttt{dig @10.30.4.10 gandal.internal AXFR} depuis une source autre que ns2 &
  Transfert refusé ; le transfert signé demandé par ns2 reste possible\\
C05 & Isolation entre tenants &
  depuis \ip{10.30.70.37} : \texttt{dig @10.30.4.12
  gw.tenant-008.\allowbreak gandal.internal A} &
  \opt{REFUSED}, sans divulgation de l'existence du nom\\
C06 & Isolation de l'infrastructure &
  depuis \ip{10.30.70.37} : \texttt{dig @10.30.4.12 xo.infra.gandal.internal A} &
  \opt{REFUSED} ; la même requête depuis le pool VPN administrateur aboutit\\
C07 & Récursion autorisée et bornée &
  depuis \ip{10.30.70.37} : \texttt{dig @10.30.4.12 example.org A} ; puis la
  même requête depuis une source hors \opt{allow-from} &
  \opt{NOERROR} dans le premier cas, absence de service dans le second\\
C08 & Autoritatif inatteignable depuis un tenant &
  depuis \ip{10.30.70.37} : \texttt{dig +time=2 +tries=1 @10.30.4.10
  gw.tenant-007.\allowbreak gandal.internal A} &
  Absence de réponse, refus journalisé sur la passerelle\\
C09 & Bascule sur perte d'une instance &
  Arrêt de ns1 puis requête sur ns2 ; arrêt de rec1 puis requête depuis un
  client dont le \opt{resolv.conf} liste \ip{.12} et \ip{.13} &
  Résolution maintenue dans les deux cas ; publication suspendue pendant
  l'arrêt de ns1\\
C10 & Idempotence de la publication &
  Émission deux fois du même appel du listing~\ref{lst:publication} &
  État final identique, aucun doublon, aucun second transfert si la lecture
  préalable est implémentée\\
C11 & Mode dégradé de publication &
  Arrêt de l'API autoritative pendant une publication &
  Opération maintenue en état observable, adresse non rendue au pool, alerte
  levée\\
C12 & Amorçage sans DNS &
  Depuis H1, accès à XO et au bastion avec le service DNS arrêté &
  Accès obtenu par le socle de reprise du listing~\ref{lst:hosts}\\
C13 & Comportement sous contrainte de MTU &
  \texttt{dig +dnssec +bufsize=1232} sur une réponse volumineuse, puis
  \texttt{dig +tcp} &
  Troncature suivie d'une reprise en TCP, sans perte silencieuse\\
C14 & Journaux et métriques en lecture pull &
  Lecture du point de métriques depuis \ip{10.30.5.10}, puis tentative
  d'initiation depuis une instance DNS vers Monitoring &
  Lecture réussie, initiation inverse refusée, horodatage en temps universel\\
C15 & Discrétion du service &
  \texttt{dig @10.30.4.10 version.bind CH TXT} & Réponse anonymisée\\
\end{longtable}

\section{Transmission à la campagne de validation}

Plusieurs scénarios du corpus dépassent le périmètre de cette tâche parce qu'ils
engagent d'autres lots ou mesurent des cibles globales. Ils sont rappelés ici
avec la contribution que T005 leur apporte, afin que la campagne de validation
puisse les reprendre sans redéfinir le contexte.

\begin{table}[H]
\centering
\small
\caption{Scénarios transmis à la campagne de validation}\label{tab:transmis}
\begin{tabular}{@{}R{0.16}R{0.34}R{0.50}@{}}
\hautfilet
\textbf{Scénario} & \textbf{Objet} & \textbf{Contribution de T005}\\
\medfilet
S04, document 03 &
  Requêtes DNS vers l'infrastructure et vers un tenant voisin depuis le VNet
  USER &
  C05 et C06 en constituent la vérification unitaire ; S04 y ajoute le parcours
  complet depuis une machine réellement provisionnée.\\
S08, document 03 &
  Restauration depuis une sauvegarde indépendante &
  La procédure et le contenu de la sauvegarde sont fournis à la
  section~\ref{sec:sauvegarde} ; la mesure du délai reste à produire.\\
G04, document 01 &
  Coupure de XO puis d'un service DNS &
  Le comportement attendu est décrit à la section~\ref{sec:suppression} :
  continuité du trafic existant, suspension des opérations nouvelles.\\
V09, document 02 &
  Création de cinquante VNets et cinq cents baux &
  La génération des cent treize zones est automatisée par le
  listing~\ref{lst:zones} ; la tenue en charge du service reste à mesurer.\\
NET-NFR-003 &
  Reconstruction en moins de soixante minutes &
  Éléments de reprise fournis aux sections~\ref{sec:amorcage}
  et~\ref{sec:sauvegarde} ; la cible est à évaluer et non à présumer.\\
\basfilet
\end{tabular}
\end{table}

% ================== Chapitre 7 : hypotheses, limites, conclusion ============
\chapter{Hypothèses, limites et points à valider}
\label{ch:limites}

La qualité d'un dossier de conception tient autant à ce qu'il établit qu'à ce
qu'il signale comme non établi. Les éléments ci-dessous distinguent les
hypothèses prises faute de donnée, les points dont la vérification conditionne
l'application, et les limites structurelles que la conception ne peut pas lever.

\section{Hypothèses formulées}

\begin{longtable}{@{}R{0.27}R{0.58}R{0.15}@{}}
\caption{Hypothèses et choix à confirmer}\label{tab:hypotheses}\\
\hautfilet
\textbf{Élément} & \textbf{Justification} & \textbf{Statut}\\
\medfilet
\endfirsthead
\hautfilet
\textbf{Élément} & \textbf{Justification} & \textbf{Statut}\\
\medfilet
\endhead
\basfilet
\endfoot
Durée de bail DHCP $B = 3\,600$ s &
  Donnée manquante, relevant de T004. Elle fixe la durée de vie des
  enregistrements dynamiques et la quarantaine par l'inégalité
  \eqref{eq:ttl}, qui reste valable si la valeur change. &
  Hypothèse\\
Quarantaine $Q = 900$ s avant réutilisation d'une adresse &
  Le document 03 laisse cette durée \statut{à valider}. La valeur proposée
  satisfait \eqref{eq:ttl} avec une marge $\delta = 300$ s. &
  Proposition\\
Récursion complète vers la racine publique &
  Préférée à une redirection vers un résolveur d'opérateur, dont le
  comportement n'est pas maîtrisé et dont la disponibilité dépendrait d'un
  équipement non manageable. & Choix\\
VNI 30404 pour le VNet Services &
  Pris dans la réserve de travail 30404 à 30408 du document 02. L'identifiant
  réellement attribué par le contrôleur doit être relevé et rapproché de cette
  réserve. & Proposition\\
Moteur de stockage SQLite &
  Suffisant pour $N_Z = 113$ zones et un régime courant $N_R \approx 1\,330$
  enregistrements. Un moteur relationnel partagé reste une extension. &
  Choix\\
\end{longtable}

\section{Points à vérifier avant application}

\begin{itemize}
  \item \textbf{Versions exactes des logiciels.} Les directives \opt{primary} et
    \opt{secondary} supposent une version 4.5 ou supérieure du serveur
    autoritatif. La branche 5 du résolveur propose un format de configuration
    YAML qu'il faut distinguer du format classique employé ici. La disponibilité
    de la journalisation dnstap dépend également des options de compilation du
    paquet retenu.
  \item \textbf{Collecteur de journalisation.} Le composant qui consomme le flux
    dnstap et produit les fichiers lus par Monitoring n'est pas arrêté. Son
    choix doit respecter la règle d'absence d'initiation vers Monitoring.
  \item \textbf{Comportement réel de la sortie externe.} La récursion vers la
    racine suppose que le routeur R1 laisse passer le port 53 sortant en UDP et
    en TCP sans altération. Le document 01 rappelle que les fonctions réelles de
    ce routeur restent à vérifier.
  \item \textbf{Correspondance entre session VPN et tenant.} Tant qu'elle n'est
    pas fournie par le lot VPN et IAM, le pool VPN utilisateur reste privé de
    résolution interne, comme expliqué à la section~\ref{sec:vpn}.
  \item \textbf{Coût du script de politique.} Le script de la
    section~\ref{sec:script} s'exécute à chaque requête. Son incidence sur la
    latence doit être mesurée à la charge attendue avant mise en service, même
    si son contenu est volontairement minimal.
\end{itemize}

\section{Limites structurelles assumées}

\begin{itemize}
  \item La redondance du service couvre la perte d'une machine virtuelle ou
    d'un hôte. Elle ne couvre ni la perte du commutateur SW1, ni celle du
    routeur R1, qui restent des dépendances communes imposées par le cadrage
    matériel.
  \item La politique de vue repose sur l'adresse source. Elle suppose donc que
    l'usurpation d'adresse à l'intérieur d'un VNet est traitée par la
    segmentation et les protections du lot VNet et SDN. Le DNS ne constitue pas,
    à lui seul, un mécanisme d'authentification du demandeur.
  \item Un tenant peut énumérer ses propres noms. Cette propriété est inhérente
    à un service de résolution ouvert à ses clients et ne franchit aucune
    frontière de tenant.
  \item Aucun déploiement n'a été réalisé et aucune mesure n'a été produite. Les
    volumes, les délais et les capacités énoncés sont des estimations de
    conception, à confronter à la maquette avant toute mise en service.
\end{itemize}

\section*{Conclusion}
\addcontentsline{toc}{section}{Conclusion}

Ce dossier fournit ce que la tâche T005 demandait : une architecture de
résolution arrêtée, un espace de nommage entièrement dérivé du plan d'adressage
du document 02, un mécanisme démontré d'isolation entre domaines, des
configurations applicables pour les quatre instances, et un cycle de vie des
enregistrements compatible avec la machine d'états du document 03.

Trois décisions en constituent l'ossature. La séparation stricte des rôles
autoritatif et récursif, qui rend l'isolation applicable plutôt que seulement
souhaitée. L'indexation de la politique de vue sur l'étiquette $\tau(a)$
calculée avant le cache, qui empêche le cache de contourner la politique. Enfin
le propriétaire d'écriture unique, qui rend la publication rejouable et la
réconciliation possible.

Le service reste conditionné aux hypothèses de ce chapitre et à la simulation
préalable exigée par \exig{net-fr-020}. L'exécution des contrôles du
chapitre~\ref{ch:recette}, la mesure des cibles de reconstruction et l'archivage
des preuves relèvent de la tâche T009, qui interviendra lorsque les lots DHCP,
temps et agent seront disponibles.


% ----------------------------------------------------------------- references
% ================== References =============================================
\cleardoublepage
\phantomsection
\addcontentsline{toc}{chapter}{Références}
\markboth{Références}{}

\begin{thebibliography}{99}
\small

\bibitem{doc01} Projet GANDAL. \textbf{Architecture globale GANDAL}, document
  01, version 1.2, 4 octobre 2026. Besoins, exigences, responsabilités et
  interfaces du datacenter.

\bibitem{doc02} Projet GANDAL. \textbf{Conception VNet et SDN GANDAL}, document
  02, version 1.2, 4 octobre 2026. Plan d'adressage détaillé, réservations VNI,
  MTU et matrice de flux. Référence des plages consommées par ce dossier.

\bibitem{doc03} Projet GANDAL. \textbf{Services réseau GANDAL}, document 03,
  version 1.2, 4 octobre 2026. IPAM, DHCP, DNS, temps et provisionnement des
  machines virtuelles. Arbitrage logiciel et machine d'états de création.

\bibitem{doc04} Projet GANDAL. \textbf{Topologie cible GANDAL}, document 04,
  version 1.2, 4 octobre 2026. Inventaire des objets à représenter et
  conventions de dessin.

\bibitem{underlay} Projet GANDAL. \textbf{Services réseau et infrastructure
  Underlay}, état de l'art technique, octobre 2026. Adressage, résolution de
  noms, dépendances de démarrage et accès d'administration.

\bibitem{suivi} Projet GANDAL. \textbf{Suivi des tâches de l'équipe 20}, 2026.
  Attribution des tâches T001 à T009.

\bibitem{rfc1034} IETF. \textbf{RFC 1034, Domain names, concepts and
  facilities}, 1987. Définition de la délégation, de l'autorité et de la mise en
  cache.

\bibitem{rfc1035} IETF. \textbf{RFC 1035, Domain names, implementation and
  specification}, 1987. Format des messages, transferts de zone et champs du
  SOA.

\bibitem{rfc2317} IETF. \textbf{RFC 2317, Classless IN-ADDR.ARPA delegation},
  1998. Découpage inverse sans classe, examiné puis écarté à la
  section~\ref{sec:nommage}.

\bibitem{rfc6303} IETF. \textbf{RFC 6303, Locally served DNS zones}, 2011.
  Fondement du comportement par défaut neutralisé par la directive
  \opt{serve-rfc1918}.

\bibitem{rfc6891} IETF. \textbf{RFC 6891, Extension mechanisms for DNS,
  EDNS(0)}, 2013. Annonce de la taille de réponse UDP, paramétrée à 1\,232
  octets.

\bibitem{rfc7766} IETF. \textbf{RFC 7766, DNS transport over TCP}, 2016.
  Justification de l'ouverture obligatoire du port 53 en TCP.

\bibitem{rfc8945} IETF. \textbf{RFC 8945, Secret key transaction authentication
  for DNS, TSIG}, 2020. Signature des transferts de zone entre ns1 et ns2.

\bibitem{rfc2136} IETF. \textbf{RFC 2136, Dynamic updates in the domain name
  system}, 1997. Mécanisme de l'alternative évaluée au
  tableau~\ref{tab:arbitrage}.

\bibitem{rfc4035} IETF. \textbf{RFC 4035, Protocol modifications for the DNS
  security extensions}, 2005. Validation DNSSEC et traitement des branches non
  signées.

\bibitem{icann} ICANN. \textbf{Board resolution on the reservation of the
  .internal top level domain for private use}, 2024. Fondement du choix de
  suffixe interne.

\bibitem{pdnsauth} PowerDNS. \textbf{Authoritative server documentation} :
  settings, HTTP API et \opt{pdnsutil}. Documentation évolutive ; la version
  installée doit être confirmée avant application.

\bibitem{pdnsrec} PowerDNS. \textbf{Recursor documentation} : settings,
  scripting Lua, \opt{gettag} et \opt{preresolve}. Documentation évolutive ; la
  version installée doit être confirmée avant application.

\bibitem{flagday} DNS Flag Day 2020. \textbf{Recommandation sur la taille des
  messages DNS en UDP}. Origine de la valeur de 1\,232 octets retenue à
  l'équation \eqref{eq:mtu}.

\bibitem{xcpng} XCP-ng Project. \textbf{Networking et XCP-ng 8.3 LTS}. Version
  cible du corpus, à confirmer sur les hôtes.

\end{thebibliography}


\end{document}
